% Chapter VII, first half: §64–71
% The Stability of Couette Flow

\chapter{THE STABILITY OF COUETTE FLOW}
\label{ch:7}

\section{Introduction}\label{sec:7-64}
In the last five chapters we have considered the onset of instability
under a variety of externally impressed conditions; but the origin of the
instability was always the same: a potentially unstable arrangement of
the fluid resulting from a prevailing adverse temperature gradient.  In this
and in the following two chapters we shall consider cases of instabilities
which have a different origin: a potentially unstable arrangement of flow
resulting from a prevailing adverse gradient of angular momentum.  The
simplest example of such instability occurs in \emph{Couette flow}, i.e.\ in the
steady circular flow of a liquid between two rotating coaxial cylinders.

In this chapter we shall consider the simple Couette flow of inviscid
and viscid fluids.  In Chapter~VIII we shall consider some examples of
more general curved flows; and in Chapter~IX we shall consider the
effect of an axial magnetic field on Couette flow.

\section{The physical problem}\label{sec:7-65}
We consider, then, the permissible stationary circular flows of an
incompressible fluid between two rotating coaxial cylinders.  As we shall
presently see, in the absence of viscosity, the class of such permissible
flows is very wide: indeed, if $\Omega$ denotes the angular velocity of rotation
about the axis, then the equations of motion allow $\Omega$ to be an arbitrary
function of the distance $r$ from the axis, provided the velocities in the
radial and the axial directions are zero.  But if viscosity is present, the
class becomes very restricted; in fact, in the absence of any transverse
pressure gradient, the most general form of $\Omega$ which is allowed is
\begin{equation}
\label{eq:7-1}
\Omega(r) = A + B/r^2,
\end{equation}
where $A$ and $B$ are two constants which are related to the angular
velocities $\Omega_1$ and $\Omega_2$ with which the inner and the outer cylinders are
rotated.  Thus, if $R_1$ and $R_2$ ($> R_1$) are the radii of the two cylinders,
then
\begin{equation}
\label{eq:7-2}
\Omega_1 = A + B/R_1^2 \quad \text{and} \quad \Omega_2 = A + B/R_2^2.
\end{equation}
Solving for $A$ and $B$ in terms of $\Omega_1$ and $\Omega_2$, we have
\begin{equation}
\label{eq:7-3}
A = -\Omega_1 \eta^2 \frac{1-\mu/\eta^2}{1-\eta^2} \quad \text{and} \quad
B = \Omega_1 \frac{R_1^2(1-\mu)}{1-\eta^2},
\end{equation}
where
\begin{equation}
\label{eq:7-4}
\mu = \Omega_2/\Omega_1 \quad \text{and} \quad \eta = R_1/R_2.
\end{equation}

The questions to which we seek answers are the following.  In the
absence of viscosity when $\Omega(r)$ can be an arbitrary function of $r$, what is
the necessary and sufficient condition for the stability of the flow?
What does this condition imply for the distribution of $\Omega$ given by~\eqref{eq:7-1}?
And, finally, how is the condition for inviscid fluids, as applied to~\eqref{eq:7-1},
modified when account is taken of viscosity?  These are the questions
to which the present chapter is addressed.

\section{Rayleigh's criterion}\label{sec:7-66}
Rayleigh stated: \emph{in the absence of viscosity, the necessary and sufficient
condition for a distribution of angular velocity $\Omega(r)$ to be stable is}
\begin{equation}
\label{eq:7-5}
\frac{d}{dr}(r^2\Omega)^2 > 0
\end{equation}
\emph{everywhere in the interval; and, further, that the distribution is unstable if
$(r^2\Omega)^2$ should decrease anywhere inside the interval.}

Since $|r^2\Omega|$ is the angular momentum, per unit mass, of a fluid element
about the axis of rotation, an alternative way of stating Rayleigh's
criterion is: \emph{a stratification of angular momentum about an axis is stable if
and only if it increases monotonically outward.}

Rayleigh did not establish his criterion by an analytical discussion of
the relevant perturbation equations.  Instead, he argued as follows.

The hydrodynamical equations governing an incompressible inviscid
fluid in cylindrical polar coordinates $(r,\theta,z)$ are:\footnote{The considerations which follow are equally applicable if an external force derivable
from an axisymmetric potential is present.}
\begin{equation}
\label{eq:7-6}
\frac{\partial u_r}{\partial t} + (\mathbf{u}.\mathrm{grad})\,u_r - \frac{u_\theta^2}{r}
= -\frac{\partial}{\partial r}\!\left(\frac{p}{\rho}\right),
\end{equation}
\begin{equation}
\label{eq:7-7}
\frac{\partial u_\theta}{\partial t} + (\mathbf{u}.\mathrm{grad})\,u_\theta + \frac{u_r u_\theta}{r}
= -\frac{1}{r}\frac{\partial}{\partial \theta}\!\left(\frac{p}{\rho}\right),
\end{equation}
\begin{equation}
\label{eq:7-8}
\frac{\partial u_z}{\partial t} + (\mathbf{u}.\mathrm{grad})\,u_z
= -\frac{\partial}{\partial z}\!\left(\frac{p}{\rho}\right),
\end{equation}
where
\begin{equation}
\label{eq:7-9}
(\mathbf{u}.\mathrm{grad}) = u_r \frac{\partial}{\partial r}
+ \frac{u_\theta}{r}\frac{\partial}{\partial \theta}
+ u_z \frac{\partial}{\partial z}.
\end{equation}
We also have the equation of continuity,
\begin{equation}
\label{eq:7-10}
\frac{\partial u_r}{\partial r}
+ \frac{u_r}{r}
+ \frac{1}{r}\frac{\partial u_\theta}{\partial \theta}
+ \frac{\partial u_z}{\partial z} = 0.
\end{equation}
Those equations clearly allow the stationary solution,
\begin{equation}
\label{eq:7-11}
u_r = u_z = 0 \quad \text{and} \quad u_\theta = V(r),
\end{equation}
where $V(r)$ is an arbitrary function of $r$.  In terms of $V(r)$, the pressure
distribution is determined by
\begin{equation}
\label{eq:7-12}
p = \rho \int \frac{dr}{r}\, V^2(r).
\end{equation}

For axisymmetric motions to which the rest of the discussion in this
section will be restricted, equations~\eqref{eq:7-6}--\eqref{eq:7-8} reduce to:
\begin{equation}
\label{eq:7-13}
\frac{\partial u_r}{\partial t}
+ u_r \frac{\partial u_r}{\partial r}
+ u_z \frac{\partial u_r}{\partial z}
- \frac{u_\theta^2}{r}
= -\frac{\partial}{\partial r}\!\left(\frac{p}{\rho}\right),
\end{equation}
\begin{equation}
\label{eq:7-14}
\frac{\partial u_\theta}{\partial t}
+ u_r \frac{\partial u_\theta}{\partial r}
+ u_z \frac{\partial u_\theta}{\partial z}
+ \frac{u_r u_\theta}{r} = 0,
\end{equation}
\begin{equation}
\label{eq:7-15}
\frac{\partial u_z}{\partial t}
+ u_r \frac{\partial u_z}{\partial r}
+ u_z \frac{\partial u_z}{\partial z}
= -\frac{\partial}{\partial z}\!\left(\frac{p}{\rho}\right).
\end{equation}
Equation~\eqref{eq:7-14} is equivalent to
\begin{equation}
\label{eq:7-16}
\frac{d}{dt}(r u_\theta) = \frac{d}{dt}(r^2\Omega) = 0.
\end{equation}
Therefore, the angular momentum $L$ ($= r^2\Omega$) of a fluid element, per unit
mass, remains constant as we follow it with its motion.  And the motions
in the radial and the axial directions take place as though $u_\theta$ were absent
and, instead, a force $u_\theta^2 r = L^2/r^3$ were acting in the radial direction.  Since
$L$ is a constant of the motion, we can associate with the force $L^2/r^3$ a
potential energy $\frac{1}{2}L^2/r^2$.  This potential energy is the same as the
kinetic energy associated with the velocity $u_\theta$ in the transverse direction.

Suppose now that we interchange the fluid contained in two elementary
rings, of equal heights and masses, at $r = r_1$ and $r = r_2$ ($> r_1$).  If
$dr_1$ and $dr_2$ are the radial extents of the rings, the equality of their
masses requires $2\pi r_1\, dr_1 = 2\pi r_2\, dr_2 = dS$ (say).  In view of the constancy
of $L$ with the motion, the fluid at $r_1$, after the interchange, will have
the same angular momentum (namely, $L_1$) which it had at $r_1$ before the
interchange; similarly, the fluid at $r_2$, after the interchange, will have the
same angular momentum (namely, $L_2$) which it had at $r_2$.  As a result,
the change in the kinetic energy (or, what is the same thing, the change
in the centrifugal potential energy) is proportional to
\begin{equation}
\label{eq:7-17}
\left(\frac{L_1^2}{r_2^2} + \frac{L_2^2}{r_1^2}\right)
- \left(\frac{L_1^2}{r_1^2} + \frac{L_2^2}{r_2^2}\right)
= (L_2^2 - L_1^2)\!\left(\frac{1}{r_1^2} - \frac{1}{r_2^2}\right)\! dS.
\end{equation}
Remembering that $r_2 > r_1$, we observe that this is positive or negative
according as $L_2^2$ is greater than $L_1^2$ or less than $L_1^2$.  Consequently, if $L^2$
is a monotonic increasing function of $r$, no interchange of fluid rings
such as we have imagined can occur without a source of energy; and this
means stability.  On the other hand, if $L^2$ should decrease anywhere,
then an interchange of fluid rings in this region will result in a liberation
of energy; and this means instability.

While the foregoing arguments of Rayleigh make his criterion a very
likely one by drawing attention to its physical origin, one should still
like to establish it directly from the relevant perturbation equations;
this we shall do in \S~67.  Meantime, it is useful to see what Rayleigh's
criterion implies for the particular distribution of $\Omega$ which is permissible
for viscid fluids.  As the \emph{discriminant} for stability, we shall use
\begin{equation}
\label{eq:7-18}
\Phi(r) = \frac{1}{r^3}\frac{d}{dr}(r^2\Omega)^2
= \frac{2\Omega}{r}\frac{d}{dr}(r^2\Omega).
\end{equation}
Rayleigh's criterion requires
\begin{equation}
\label{eq:7-19}
\Phi(r) > 0 \quad \text{for stability.}
\end{equation}
When $\Omega$ has the form given by~\eqref{eq:7-1},
\begin{equation}
\label{eq:7-20}
\Phi = 4A(A + B/r^4).
\end{equation}
It is convenient to measure $r$ in units of the radius $R_2$ of the outer
cylinder.  Then
\begin{equation}
\label{eq:7-21}
\frac{4AB}{R_2^4}\!\left(\frac{1}{r^4}+\frac{A\,R_2^4}{B}\right).
\end{equation}
With the values of $A$ and $B$ given in~\eqref{eq:7-3}, we have
\begin{equation}
\label{eq:7-22}
\Phi = -4\Omega_1^2\,\eta^4\,\frac{(1-\mu)(1-\mu/\eta^2)}{(1-\eta^2)^2}
\!\left(\frac{1}{r^4} - \kappa\right),
\end{equation}
where
\begin{equation}
\label{eq:7-23}
\kappa = \frac{A\, R_2^4}{B} = \frac{1-\mu/\eta^2}{1-\mu}.
\end{equation}
Clearly,
\begin{equation}
\label{eq:7-24}
\Phi > 0 \quad \text{for} \quad \eta \leqslant r \leqslant 1,
\end{equation}
so long as
\begin{equation}
\label{eq:7-25}
\mu = \Omega_2/\Omega_1 > \eta^2.
\end{equation}
Therefore, Rayleigh's criterion applied to the distribution~\eqref{eq:7-1} requires
that \emph{for stability, the outer cylinder must rotate with an angular speed
greater than $\eta^2$ times that of the inner cylinder and in the same sense.}  In
the $(\Omega_1, \Omega_2)$-plane\footnote{Without loss of generality, we may suppose that the sense of rotation of the inner
cylinder is positive.} (see Fig.~62), the regions of stability are delimited
by the positive $\Omega_2$-axis and the Rayleigh line $\Omega_2 = \Omega_1 \eta^2$.  If the effects
of viscosity are ignored, we must have instability to the left of the
Rayleigh line; and it is of interest to determine how Rayleigh's criterion
is violated for flow represented in this part of the plane.  With $\kappa$ given
by~\eqref{eq:7-23}, we find
\begin{equation}
\label{eq:7-26}
\frac{1}{r^4} - \kappa = \begin{cases}
\mu\,\dfrac{1-\eta^2}{\eta^4(1-\mu)} & \text{for } r = 1,\\[6pt]
\dfrac{(1-\eta^2)}{\eta^4(1-\mu)} & \text{for } r = \eta.
\end{cases}
\end{equation}
Accordingly,
\begin{equation}
\label{eq:7-27}
\Phi(r) < 0 \quad \text{for} \quad \eta \leqslant r \leqslant 1,
\end{equation}
so long as
\begin{equation}
\label{eq:7-28}
0 < \mu < \eta^2.
\end{equation}
Therefore, for $0 < \mu < \eta^2$, Rayleigh's criterion is violated everywhere

\figurePlaceholder{62}{The $(\Omega_1,\Omega_2)$-plane: in the limit of zero viscosity,
the stable and the unstable regions in this plane are separated
by the Rayleigh line, $\Omega_2 = \Omega_1(R_1/R_2)^2$; $R_1$ and $R_2$ are the
radii and $\Omega_1$ and $\Omega_2$ are the angular velocities of the inner
and the outer cylinders, respectively.}

in the fluid; and we may say that instability occurs throughout the
fluid.  On the other hand, if the cylinders are rotating in opposite
directions and $\mu < 0$, $\Phi$ is positive in a part of the fluid adjoining the
outer cylinder; and it is negative in the remaining part of the fluid
adjoining the inner cylinder.  The two parts are separated by the \emph{nodal
surface} on which $\Omega = 0$.  The radius $r_0$ of this surface in the unit $R_2$ is
given by
\begin{equation}
\label{eq:7-29}
r_0 = \frac{1}{\sqrt{\kappa}} = \eta\sqrt{\frac{1+|\mu|}{{\eta^2}+|\mu|}}\,;
\end{equation}
and
\begin{equation}
\label{eq:7-30}
\left.\begin{aligned}
\Phi(r) &> 0 \quad \text{for} \quad r_0 < r < 1\\
\Phi(r) &< 0 \quad \text{for} \quad \eta < r < r_0
\end{aligned}\right\}.
\end{equation}
Thus, when $\mu < 0$, Rayleigh's criterion is violated only interior to the
nodal surface.  The instability which is predicted for all negative $\mu$ in
the absence of viscosity, is therefore, to be associated with this instability
of the flow interior to $r_0$.  On this account we may say that for $\mu < 0$,
the instability is only partial.  Also, as $\mu \to -\infty$, $r_0 \to \eta$ and the part
of the fluid which is unstable according to Rayleigh's criterion becomes
vanishingly small.

If we now consider the effect of viscosity on the stability of the flow
prescribed by~\eqref{eq:7-1}, we must on general grounds expect that it will postpone the onset of instability beyond the point predicted by Rayleigh's
criterion; and that for a given $\kappa$ (i.e.\ $\mu$ and $\eta$) the factor,
\begin{equation}
\label{eq:7-31}
4\Omega_1^2\,\eta^4\,\frac{(1-\mu)(1-\mu/\eta^2)}{(1-\eta^2)^2}
\quad (\mu < \eta^2),
\end{equation}
which occurs in Rayleigh's discriminant must exceed a certain critical
value, depending on $\nu$ and $R_2$, before instability can set in.  The non-dimensional number in terms of which the criterion for stability will be
expressed is, then,
\begin{equation}
\label{eq:7-32}
T = \frac{4\Omega_1^2}{\nu^2}\,R_2^4\,\frac{(1-\mu)(1-\mu/\eta^2)}{(1-\eta^2)^2}.
\end{equation}
This is the proper definition of the \emph{Taylor number} for this problem.  For
a given $\eta$ and $\mu$, instability will set in for a certain critical Taylor number
$T_c$; and the central problem of this subject is the determination of $T_c$ as a
function of $\eta$ and $\mu$.

\section{Analytical discussion of the stability of inviscid Couette flow}
\label{sec:7-67}
We now return to an analytical discussion of the stability of the
stationary flow,
\begin{equation}
\label{eq:7-33}
u_r = u_z = 0 \quad \text{and} \quad u_\theta = V(r) = r\Omega(r),
\end{equation}
which the equations of inviscid motion~\eqref{eq:7-6}--\eqref{eq:7-9} allow.
In~\eqref{eq:7-33}, $V(r)$ is an arbitrary function of $r$.

Consider an infinitesimal perturbation of the flow represented by the
solution~\eqref{eq:7-33}.  Let the perturbed state be described by
\begin{equation}
\label{eq:7-34}
u_r, \quad V\!+\!u_\theta, \quad u_z, \quad \text{and} \quad
\varpi\;(= \delta p/\rho).
\end{equation}
The linear equations governing these perturbations are
\begin{equation}
\label{eq:7-35}
\frac{\partial u_r}{\partial t}
+ \frac{V}{r}\frac{\partial u_r}{\partial \theta}
- 2\frac{V}{r}\,u_\theta
= -\frac{\partial \varpi}{\partial r},
\end{equation}
\begin{equation}
\label{eq:7-36}
\frac{\partial u_\theta}{\partial t}
+ \frac{V}{r}\frac{\partial u_\theta}{\partial \theta}
+ \left(\frac{dV}{dr}+\frac{V}{r}\right) u_r
= -\frac{1}{r}\frac{\partial \varpi}{\partial \theta},
\end{equation}
\begin{equation}
\label{eq:7-37}
\frac{\partial u_z}{\partial t}
+ \frac{V}{r}\frac{\partial u_z}{\partial \theta}
= -\frac{\partial \varpi}{\partial z},
\end{equation}
and the equation of continuity,
\begin{equation}
\label{eq:7-38}
\frac{\partial u_r}{\partial r}
+ \frac{u_r}{r}
+ \frac{1}{r}\frac{\partial u_\theta}{\partial \theta}
+ \frac{\partial u_z}{\partial z} = 0.
\end{equation}

In accordance with the general procedure of treating these problems,
we analyse the disturbance into normal modes.  In the present instance,
it is natural to suppose that the various quantities describing the
perturbation have a $(t,\theta,z)$-dependence given by
\begin{equation}
\label{eq:7-39}
e^{i(\sigma t + m\theta + kz)},
\end{equation}
where $p$ is a constant (which can be complex), $m$ is an integer (which
can be positive, zero, or negative), and $k$ is the wave number of the
disturbance in the $z$-direction.

Let $u_r(r)$, $u_\theta(r)$, $u_z(r)$, and $\varpi(r)$ now denote the amplitudes of the
respective perturbations whose $(t,\theta,z)$-dependence is given by~\eqref{eq:7-39}.
Equations~\eqref{eq:7-35}--\eqref{eq:7-38} then give
\begin{equation}
\label{eq:7-40}
i\sigma u_r - 2\Omega u_\theta = -\frac{d\varpi}{dr},
\end{equation}
\begin{equation}
\label{eq:7-41}
i\sigma u_\theta + \left(\Omega + \frac{d}{dr} r\Omega\right) u_r
= -\frac{im}{r}\,\varpi,
\end{equation}
\begin{equation}
\label{eq:7-42}
i\sigma u_z = -ik\varpi,
\end{equation}
and
\begin{equation}
\label{eq:7-43}
\frac{du_r}{dr} + \frac{u_r}{r} + \frac{im}{r}\, u_\theta + ik u_z = 0,
\end{equation}
where
\begin{equation}
\label{eq:7-44}
\sigma = p + m\Omega.
\end{equation}
(Note that through $\Omega$, $\sigma$ is a function of $r$.)

\subsection*{(a) \emph{The equations in terms of the Lagrangian displacement}}
Consider the variables $\xi_r$, $\xi_\theta$, and $\xi_z$ related to $u_r$, $u_\theta$, and $u_z$ by
\begin{equation}
\label{eq:7-45}
u_r = i\sigma \xi_r, \quad
u_\theta = i\sigma \xi_\theta - r\frac{d\Omega}{dr}\,\xi_r, \quad \text{and} \quad
u_z = i\sigma \xi_z.
\end{equation}
Defined in this manner, $\boldsymbol{\xi}$ is the \emph{Lagrangian displacement} appropriate to
the problem on hand.

It can be readily verified that consistent with the meaning of $\boldsymbol{\xi}$ as the
Lagrangian displacement, the solenoidal character of $\mathbf{u}$ implies the
solenoidal character of $\boldsymbol{\xi}$; thus
\begin{equation}
\label{eq:7-46}
\frac{d\xi_r}{dr} + \frac{\xi_r}{r} + \frac{im}{r}\,\xi_\theta + ik\xi_z = 0.
\end{equation}
In terms of the variables $\xi_r$, $\xi_\theta$, and $\xi_z$, equations~\eqref{eq:7-40}--\eqref{eq:7-42} become
\begin{equation}
\label{eq:7-47}
\left(\sigma^2 - 2r\Omega\frac{d\Omega}{dr}\right)\xi_r + 2i\Omega\sigma\xi_\theta
= \frac{d\varpi}{dr},
\end{equation}
\begin{equation}
\label{eq:7-48}
\sigma^2 \xi_\theta - 2i\Omega\sigma\xi_r = \frac{im}{r}\,\varpi,
\end{equation}
and
\begin{equation}
\label{eq:7-49}
\sigma^2 \xi_z = ik\varpi.
\end{equation}
Multiplying equations~\eqref{eq:7-48} and~\eqref{eq:7-49} by $im/r$ and $ik$, respectively,
adding, and making use of equation~\eqref{eq:7-46}, we obtain
\begin{equation}
\label{eq:7-50}
\sigma^2\!\left(\frac{d\xi_r}{dr}+\frac{\xi_r}{r}\right)
+ 2\frac{m\Omega}{r}\sigma\xi_r
= \left(\frac{m^2}{r^2}+k^2\right)\varpi.
\end{equation}
Next eliminating $\xi_\theta$ between equations~\eqref{eq:7-47} and~\eqref{eq:7-48}, we have
\begin{equation}
\label{eq:7-51}
\left(\sigma^2 - 2r\Omega\frac{d\Omega}{dr}\right)\xi_r
+ \frac{2\Omega}{\sigma}\!\left(2i\Omega\sigma\xi_r
+ \frac{im}{r}\,\varpi\right)
= \frac{d\varpi}{dr}.
\end{equation}
Rearranging the terms in this equation, we obtain
\begin{equation}
\label{eq:7-52}
[\sigma^2 - \Phi(r)]\xi_r = \frac{d\varpi}{dr}
+ \frac{2m\Omega}{r\sigma}\,\varpi,
\end{equation}
where
\begin{equation}
\label{eq:7-53}
\Phi(r) = 2r\Omega\frac{d\Omega}{dr} + 4\Omega^2
= \frac{2\Omega}{r}\frac{d}{dr}(r^2\Omega)
\end{equation}
is the Rayleigh discriminant we have already introduced in \S~66 (equation~\eqref{eq:7-18}).

Equation~\eqref{eq:7-52} must be considered in conjunction with equation~\eqref{eq:7-50}, which we shall rewrite in the form
\begin{equation}
\label{eq:7-54}
\frac{1}{r}\frac{d}{dr}(r\xi_r)
- \frac{2m\Omega}{\sigma r}\,\xi_r
= \frac{1}{\sigma^2}\!\left(\frac{m^2}{r^2}+k^2\right)\varpi.
\end{equation}

If the fluid is confined between two coaxial cylinders of radii $R_1$ and
$R_2$, we must require that the radial components of the velocity vanish
for these values of $r$.  Thus, equations~\eqref{eq:7-52} and~\eqref{eq:7-54} must be considered
together with the boundary conditions
\begin{equation}
\label{eq:7-55}
\xi_r = 0 \quad \text{for } r = R_1 \text{ and } R_2.
\end{equation}

We shall presently see that equations~\eqref{eq:7-52} and~\eqref{eq:7-54}, for real $\sigma$, constitute a self-adjoint system with respect to the boundary conditions~\eqref{eq:7-55}; and that in consequence their solution can be reduced to a variational problem; this was, indeed, the object of introducing the Lagrangian
displacement as the variable.

\subsection*{(b) \emph{The case $m = 0$}}
When $m = 0$, $\sigma = p$ and equations~\eqref{eq:7-52} and~\eqref{eq:7-54} become
\begin{equation}
\label{eq:7-56}
[p^2 - \Phi(r)]\xi_r = \frac{d\varpi}{dr},
\end{equation}
and
\begin{equation}
\label{eq:7-57}
\frac{1}{r}\frac{d}{dr}(r\xi_r) = \frac{k^2}{p^2}\,\varpi.
\end{equation}
Eliminating $\varpi$ between these equations, we obtain
\begin{equation}
\label{eq:7-58}
\frac{d}{dr}\!\left\{
\frac{1}{r}\frac{d}{dr}(r\xi_r)\right\}
- k^2\xi_r = -\frac{k^2}{p^2}\,\Phi(r)\xi_r.
\end{equation}
The solution of this equation, together with the boundary conditions~\eqref{eq:7-55},
constitutes a characteristic value problem of the classical Sturm--Liouville type; and by appealing to the standard theorems of the subject,
we can reach the following conclusion.

\emph{The characteristic values of $k^2\!/p^2$ are all positive if $\Phi(r)$ is everywhere
positive; and they are all negative if $\Phi(r)$ is everywhere negative.  If $\Phi(r)$
should change sign anywhere in the interval $(R_1, R_2)$, then there are two
sets of real characteristic values which have the limit points $+\infty$ and $-\infty$.}

Since a negative $p^2$ means instability, it is clear that the foregoing
statements regarding the sign of the characteristic values of $k^2\!/p^2$ are
equivalent to a restatement of Rayleigh's criterion.

An alternative way of establishing the same results is instructive.
Multiply equation~\eqref{eq:7-58} by $r\xi_r$ and integrate over the range of $r$.  The
integral on the left-hand side can be transformed by an integration by
parts.  We obtain
\begin{equation}
\label{eq:7-59}
\int\!\left\{
\frac{1}{r}\!\left(\frac{d}{dr}r\xi_r\right)^{\!2}
+ k^2 r\xi_r^2
\right\} dr
= \frac{k^2}{p^2}\int \Phi(r)\, r\,\xi_r^2\, dr
\end{equation}
or
\begin{equation}
\label{eq:7-60}
\frac{p^2}{k^2} = \frac{\displaystyle\int \Phi(r)\, r\, \xi_r^2\, dr}
{\displaystyle\int\!\left\{
\frac{1}{r}[\!(d/dr)\, r\xi_r\!]^2 + k^2 r\xi_r^2
\right\} dr}.
\end{equation}
From equation~\eqref{eq:7-60} it is immediately apparent that $p^2\!/k^2$ is positive if
$\Phi(r)$ is everywhere positive, and is negative if $\Phi(r)$ is everywhere negative.
The further result that there exist unstable modes, if $\Phi(r)$ is anywhere
negative, can be deduced from the fact that we may regard equation~\eqref{eq:7-58} as the condition that
\begin{equation}
\label{eq:7-61}
I_1 = \int \Phi(r)\, r\, \xi_r^2\, dr
\end{equation}
is a maximum, or a minimum, for a given
\begin{equation}
\label{eq:7-62}
I_2 = \int\!\left\{
\frac{1}{r}\!\left(\frac{d}{dr}r\xi_r\right)^{\!2}
+ k^2 r\xi_r^2
\right\} dr.
\end{equation}
$p^2\!/k^2$ being, then, the value of $I_1/I_2$.  If $\Phi(r)$ should anywhere be negative,
$I_1$ admits a negative value and therefore a negative minimum, so that
one value at least of $p^2$ is negative, and one mode of disturbance is
unstable.

By comparing this alternative method of establishing Rayleigh's
criterion with the corresponding method of establishing the stability of
an incompressible fluid of variable density (Chapter~X, \S\,92), we observe
that $\Phi(r)$ in this problem, and the density gradient in the other problem,
play entirely equivalent roles.

\subsection*{(c) \emph{The case $m \neq 0$}}
The treatment of the general case is not as straightforward as the
case $m = 0$.  It appears that when $m \neq 0$, it is more convenient to consider
the problem of solving equations~\eqref{eq:7-52} and~\eqref{eq:7-54} subject to the boundary
conditions~\eqref{eq:7-55} as a characteristic value problem for $k^2$ for a given $p$,
rather than as a problem for $p$ for a given $k^2$.  This inversion of the usual
procedure is advantageous in so far as the problem can then be formulated
in terms of a variational principle.  Moreover, we shall find that for real
$p$ the problem is Hermitian and the characteristic values of $k^2$ are real.
The question now arises whether the solution of the stability problem
can be inferred from the solution of the characteristic value problem
inverted in the manner we have proposed.  That it should, in principle,
be possible to infer the solution becomes clear when we realize that there
is a unique dispersion relation which gives $p$ as an analytic function of $k^2$
(with, possibly, more than one branch); and that the relationship is independent of how it is determined.  And so long as the relation between
$p$ and $k^2$ is analytic, we should be able to answer the basic questions
concerning stability by determining the relation between real $p$ and $k^2$
(positive or negative) instead of, as is more usual, by determining the
relation between positive $k^2$ and $p$ (real or complex).

Returning to equations~\eqref{eq:7-52}--\eqref{eq:7-55}, we shall first show that for real $p$
the characteristic values of $k^2$ are real.  To show this, we shall multiply
equation~\eqref{eq:7-52} by $r\xi_r^*$ and integrate over the range of $r$.  We obtain
\begin{equation}
\label{eq:7-63}
\int r[\sigma^2-\Phi(r)]|\xi_r|^2\,dr
= \int \xi_r^*\!\left(r\frac{d\varpi}{dr}+\frac{2m\Omega}{\sigma}\,\varpi\right)dr.
\end{equation}
After an integration by parts, the right-hand side becomes
\begin{equation}
\label{eq:7-64}
-\int \varpi\!\left(\frac{d}{dr}(r\xi_r^*) - \frac{2m\Omega}{\sigma}\,\xi_r^*\right) dr;
\end{equation}
and making use of equation~\eqref{eq:7-54}, we can finally write
\begin{equation}
\label{eq:7-65}
\int r[\sigma^2-\Phi(r)]|\xi_r|^2\,dr
= -\int \frac{1}{\sigma^2}
\!\left(\frac{m^2}{r}+k^2 r\right)|\varpi|^2\,dr.
\end{equation}
By replacing every term in this equation by its complex conjugate, we
infer the identity of $k^2$ and $k^{*2}$; $k^2$ is, therefore, real.  When $k^2$ is real, the
proper functions, $\xi_r$ and $\varpi$, belonging to a particular $k^2$ (and a real $p$)
are also real, and we can rewrite equation~\eqref{eq:7-65} in the manner
\begin{equation}
\label{eq:7-66}
k^2 = \frac{\displaystyle\int
r[(\Phi-\sigma^2)\xi_r^2 - (m^2\varpi^2)/(r^2\sigma^2)]\,dr}
{\displaystyle\int r\,\varpi^2/\sigma^2\,dr}
\quad (\text{say}).
\end{equation}

We shall now show that the expression~\eqref{eq:7-66} for $k^2$ provides the basis
for a variational formulation of the problem.

First, it should be remarked that in evaluating $k^2$ in accordance with~\eqref{eq:7-66}, we must suppose that $\xi_r$ is determined in terms of $\varpi$ by the equation
\begin{equation}
\label{eq:7-67}
(\sigma^2-\Phi)\xi_r = \frac{d\varpi}{dr}+\frac{2m\Omega}{r\sigma}\,\varpi;
\end{equation}
and that an allowed $\varpi$ accordingly satisfies the boundary conditions
\begin{equation}
\label{eq:7-68}
\frac{d\varpi}{dr}+\frac{2m\Omega}{\sigma}\,\varpi = 0
\quad \text{for } r = R_1 \text{ and } R_2.
\end{equation}

Consider now the effect on $k^2$ of an arbitrary variation $\delta\varpi$ in $\varpi$ compatible only with the boundary conditions~\eqref{eq:7-68}.  If $\delta\xi_r$ denotes the corresponding variation in $\xi_r$, then, by~\eqref{eq:7-67}
\begin{equation}
\label{eq:7-69}
(\sigma^2-\Phi)\,\delta\xi_r = \frac{d}{dr}\,\delta\varpi
+ \frac{2m\Omega}{\sigma r}\,\delta\varpi;
\end{equation}
and
\begin{equation}
\label{eq:7-70}
\delta\xi_r = 0 \quad \text{for } r = R_1 \text{ and } R_2.
\end{equation}
From equation~\eqref{eq:7-66} it follows that
\begin{equation}
\label{eq:7-71}
\delta k^2 = \frac{1}{I_2}\bigl(\delta I_1 - \frac{I_1}{I_2}\,\delta I_2\bigr)
= \frac{1}{I_2}\bigl(\delta I_1 - k^2 \delta I_2\bigr),
\end{equation}
where
\begin{equation}
\label{eq:7-72}
\delta I_1 = 2\int r\bigl[(\Phi-\sigma^2)\,\xi_r\,\delta\xi_r
- \frac{m^2\varpi}{r^2\sigma^2}\,\delta\varpi\bigr]\,dr
\end{equation}
and
\begin{equation}
\label{eq:7-73}
\delta I_2 = 2\int r\varpi(\delta\varpi/\sigma^2)\,dr
\end{equation}
are the variations in $I_1$ and $I_2$ resulting from the variation in $\varpi$.

The terms in $\delta\xi_r$ in the expression~\eqref{eq:7-72} for $\delta I_1$ are the same as
\begin{equation}
\label{eq:7-74}
-2\int r\xi_r\!\left(\frac{d}{dr}\,\delta\varpi+\frac{2m\Omega}{\sigma r}\,\delta\varpi\right)dr.
\end{equation}
After an integration by parts,~\eqref{eq:7-74} can be rewritten in the form
\begin{equation}
\label{eq:7-75}
2\int \delta\varpi\!\left(\frac{d}{dr}(r\xi_r)
- \frac{2m\Omega}{\sigma}\,\xi_r\right) dr.
\end{equation}
The first-order change in $k^2$ resulting from the variation in $\varpi$ is, therefore,
\begin{equation}
\label{eq:7-76}
\delta k^2 = \frac{2}{I_2}\int \delta\varpi\!\left(
\frac{1}{r}\frac{d}{dr}(r\xi_r)
- \frac{2m\Omega}{\sigma r}\,\xi_r
- \frac{1}{\sigma^2}\!\left(\frac{m^2}{r^2}+k^2\right)\varpi
\right)r\,dr.
\end{equation}
Consequently, $\delta k^2 = 0$ for an arbitrary variation $\delta\varpi$ compatible with
the boundary conditions~\eqref{eq:7-68} if and only if
\begin{equation}
\label{eq:7-77}
\frac{1}{r}\frac{d}{dr}(r\xi_r)
- \frac{2m\Omega}{\sigma r}\,\xi_r
= \frac{1}{\sigma^2}\!\left(\frac{m^2}{r^2}+k^2\right)\varpi.
\end{equation}
Conversely, whenever equation~\eqref{eq:7-77} is satisfied, $\delta k^2 = 0$.  But equation~\eqref{eq:7-77} is the same as equation~\eqref{eq:7-54}.  Therefore, solving the characteristic
value problem presented by equations~\eqref{eq:7-52}, \eqref{eq:7-54}, and~\eqref{eq:7-55} is equivalent
to finding a maximum or a minimum of
\begin{equation}
\label{eq:7-78}
I_1 = -\int \frac{1}{r^2 + \frac{m^2}{r^2}}\!\left[
\!\left(\frac{d\varpi}{dr}+\frac{2m\Omega}{\sigma r}\,\varpi\right)^{\!2}
+ \frac{m^2\varpi^2}{r^2\sigma^2}\right]\,dr,
\end{equation}
for given
\begin{equation}
\label{eq:7-79}
I_2 = \int dr\; r\varpi^2/\sigma^2,
\end{equation}
for arbitrary variations of $\varpi$ subject only to the boundary conditions
\begin{equation}
\label{eq:7-80}
\frac{d\varpi}{dr}+\frac{2m\Omega}{\sigma}\,\varpi = 0
\quad \text{for } r = R_1 \text{ and } R_2;
\end{equation}
and the ratio $I_1/I_2$ at such a maximum or a minimum is a characteristic
value of $k^2$.

From equation~\eqref{eq:7-66} it is at once apparent that if $\Phi(r)$ is everywhere
negative,\footnote{As is the case when the flow is that permissible for viscid fluids and $0<\mu<\eta^2$
(cf.\ equations~\eqref{eq:7-27a} and~\eqref{eq:7-28}).}
$k^2$ cannot admit a positive characteristic value.  For real $p$'s,
the characteristic values of $k$ are necessarily imaginary.  Therefore
for real $k$'s, $p$ must necessarily be complex; and this means instability.

Consider next the case when $\Phi(r)$ changes sign in the interval $(R_1, R_2)$
and it is negative over a part of the interval and positive over others.
Under these circumstances, $I_1$ can assume a negative value for any
arbitrarily assigned real $p$; and for a given $I_2$, $I_1$ can attain a negative
minimum.  Therefore, for any real $p$, there exists at least one negative
characteristic value for $k^2$.  At the same time, since $\Phi(r)$ is positive over
some other parts of the interval $I_1$ can also assume positive values for
suitably chosen real values of $p$.  For a given $I_2$, there exist values of
$p$ for which $I_1$ can attain a positive maximum.  Therefore, for some
real values of $p$, there exist positive characteristic values for $k^2$.  From
our earlier arguments, these same values of $p$ must also allow negative
characteristic values of $k^2$.  Thus, in this case the dispersion relation
has two branches.  One of them is such that for all real $p$'s the characteristic values of $k^2$ are negative.  The inverse of this latter branch must
lead to complex $p$'s for real $k$'s; and this means instability.

Consider finally the case when $\Phi(r)$ is everywhere positive.  In this
case it is convenient to rewrite equation~\eqref{eq:7-66} in the form
\begin{equation}
\label{eq:7-81}
k^2 = \frac{\displaystyle\int
r\{[(d\varpi/dr)+(2m\Omega/\sigma)\varpi]^2(\Phi-\sigma^2)
- (m^2\varpi^2/r^2\sigma^2)\}\,dr}
{\displaystyle\int r\,\varpi^2/\sigma^2\,dr}.
\end{equation}
It is clear that for all real $p$ such that $\sigma^2 = (p+m\Omega)^2 > \Phi(r)$, the characteristic values of $k^2$ are negative.  For positive characteristic values it is
clearly necessary that $\sigma^2$ be less than $\Phi(r)$ at least over a part of the
interval $(R_1, R_2)$.  Indeed, it is clear that by requiring $\sigma^2$ to be less than
$\Phi(r)$ over a part of the interval, we can obtain for $k^2$ arbitrary positive
values; i.e.\ for all real $k$'s there correspond real $p$'s.  The system is
therefore stable.  The necessary and sufficient character of Rayleigh's
criterion for the stability of inviscid rotational flow is therefore established for arbitrary perturbations.

\section{The periods of oscillation of a rotating column of liquid}
\label{sec:7-68}
Rayleigh's criterion, for the stability of a given distribution of angular
velocity in a cylindrical column of liquid rotating about its axis, was
established in \S~67 from a consideration of the characteristic value problem which relates the frequency of oscillation with the wavelength of
the disturbance in the axial direction.  It is of interest to obtain the
explicit form of this relation in some special cases.  As Lord Kelvin has
remarked in his first investigation of the subject, `Crowds of extremely
interesting cases present themselves'.

\subsection*{(a) \emph{The case $\Omega = \text{constant}$}}
When $\Omega = \text{constant}$,
\begin{equation}
\label{eq:7-82}
\Phi(r) = \frac{1}{r^3}\frac{d}{dr}(r^2\Omega)^2 = 4\Omega^2,
\end{equation}
and the relevant equations are (cf.\ equations~\eqref{eq:7-52} and~\eqref{eq:7-54}):
\begin{equation}
\label{eq:7-83}
(\sigma^2 - 4\Omega^2)\xi_r = \frac{d\varpi}{dr}
+ \frac{2m\Omega}{\sigma r}\,\varpi,
\end{equation}
and
\begin{equation}
\label{eq:7-84}
\frac{1}{r}\frac{d}{dr}(r\xi_r)
- \frac{2m\Omega}{\sigma}\,\xi_r
= \frac{1}{\sigma^2}\!\left(\frac{m^2}{r^2}+k^2\right)\varpi,
\end{equation}
where
\begin{equation}
\label{eq:7-85}
\sigma = p + m\Omega
\end{equation}
is now a constant.  Eliminating $\xi_r$ between equations~\eqref{eq:7-83} and~\eqref{eq:7-84}, we
obtain
\begin{equation}
\label{eq:7-86}
\frac{d^2\varpi}{dr^2}+\frac{1}{r}\frac{d\varpi}{dr}
- \frac{m^2}{r^2}\,\varpi
= -k^2\!\left(\frac{4\Omega^2}{\sigma^2}-1\right)\varpi.
\end{equation}
It is convenient to measure $r$ in units of the radius $R_2$ of the outer
cylinder and let
\begin{equation}
\label{eq:7-87}
a = kR_2 \quad \text{and} \quad R_1 = R_2\eta.
\end{equation}
Equation~\eqref{eq:7-86} then becomes
\begin{equation}
\label{eq:7-88}
\frac{d^2\varpi}{dr^2}+\frac{1}{r}\frac{d\varpi}{dr}
+ \left(a^2\!\left(\frac{4\Omega^2}{\sigma^2}-1\right)
- \frac{m^2}{r^2}\right)\varpi = 0;
\end{equation}
and the boundary conditions are (cf.\ equation~\eqref{eq:7-68})
\begin{equation}
\label{eq:7-89}
\frac{d\varpi}{dr}+\frac{2m\Omega}{\sigma}\,\varpi = 0
\quad \text{for } r = 1 \text{ and } \eta.
\end{equation}
The general solution of equation~\eqref{eq:7-88} is
\begin{equation}
\label{eq:7-90}
\varpi = A J_m(\alpha r) + B Y_m(\alpha r),
\end{equation}
where
\begin{equation}
\label{eq:7-91}
\alpha = a(4\Omega^2/\sigma^2-1)^{1/2},
\end{equation}
$A$ and $B$ are constants of integration, and $J_m$ and $Y_m$ are the Bessel functions of the two kinds of order $m$.  The application of the boundary
conditions~\eqref{eq:7-89} to the solution~\eqref{eq:7-90} leads to a transcendental equation
relating $\alpha$ and $\sigma$ ($= p+m\Omega$).  We shall consider two special cases when
the relation is fairly simple.

\subsubsection*{(i) \emph{The case $\eta = 0$}}
When there is no inner cylinder and $\eta = 0$, the freedom from
singularity at $r = 0$ requires that the term in $Y_m$ in the solution~\eqref{eq:7-90} be
absent.  We then have
\begin{equation}
\label{eq:7-92}
\varpi = A J_m(\alpha r);
\end{equation}
and the boundary condition at $r = 1$ gives
\begin{equation}
\label{eq:7-93}
\alpha J_m'(\alpha) + \frac{2m\Omega}{\sigma}\, J_m(\alpha) = 0,
\end{equation}
where a prime denotes differentiation with respect to the argument.
But according to equation~\eqref{eq:7-91},
\begin{equation}
\label{eq:7-94}
\frac{\sigma}{2\Omega} = \frac{\pm 1}{\sqrt{1+\alpha^2/a^2}}.
\end{equation}
Using this relation to eliminate $\Omega/\sigma$ in equation~\eqref{eq:7-93}, we obtain
\begin{equation}
\label{eq:7-95}
\alpha J_m'(\alpha) \pm m(1+\alpha^2/a^2)^{1/2} J_m(\alpha) = 0.
\end{equation}
For any assigned value of $\alpha$ and $m$, equation~\eqref{eq:7-95} determines $\alpha/a$; equation~\eqref{eq:7-94} then determines $p$; thus
\begin{equation}
\label{eq:7-96}
\frac{p}{\Omega} = \frac{\sigma}{\Omega} - m = \pm\frac{2}{\sqrt{1+\alpha^2/a^2}} - m.
\end{equation}
In this way the dispersion relation between $\alpha$ and $p$ can be established.

The case $m = 0$ is particularly simple.  In this case $\alpha$ must be a zero
of $J_0'(x)$, i.e.\ of $J_1(x)$.  If $\alpha_{0,j}$ denotes the $j$th zero of $J_1(x)$, the explicit
form of the dispersion relation is

\figurePlaceholder{63}{The characteristic frequencies of oscillation of a rotating column of
fluid of radius $R_2$ for various values of $m$ and $j$; $a$ measures the wave number
in the unit $1/R_2$.}

\begin{equation}
\label{eq:7-97}
p = \pm\frac{2\Omega}{\sqrt{1+\alpha_{0,j}^2/a^2}}.
\end{equation}
The dispersion relations for $m = 0$, 1, 2, and 3 are illustrated in Fig.~63.
If waves of the type we have been considering should occur as standing
waves in a rotating column of liquid, then their half-wavelengths should
divide the height $H$ of the column an integral number of times; and we
must have
\begin{equation}
\label{eq:7-98}
a = n\pi R_2/H,
\end{equation}
where $n$ is an integer.  By inserting this value of $a$ in equations~\eqref{eq:7-96} and~\eqref{eq:7-97}, we shall obtain the frequencies of the various natural modes of
oscillation of the column.

In some recent experiments, Fultz has shown how these natural modes
of oscillation of a rotating column of liquid can be excited and maintained.  In these experiments the modes $m = 0$ in a rotating cylinder of
water (of diameter 4.2~cm and height 15.3~cm) were excited by means of
a small disk on the axis of the cylinder.  This disk could be moved up
and down the axis with a frequency which could be controlled.  To excite
a particular mode, the disk was so placed that its mean height was at the
position of the topmost antinode of the mode (e.g.\ at the middle of the
column to excite the fundamental mode).  Fultz was able to detect with
great sensitivity whether the disk was set at the resonant frequency by
following the motions of the fluid by means of dye suitably introduced.
In this way Fultz determined the frequencies of some of the lower modes
of oscillation belonging to $m = 0$; and he found that the observed
resonant frequencies agreed remarkably well with those given by
equations~\eqref{eq:7-97} and~\eqref{eq:7-98}.  In Fig.~64 we reproduce some of Fultz's photographs which exhibit in a very striking manner these natural modes of
oscillation of a rotating fluid.

\figurePlaceholder{64}{Photographs illustrating the oscillations of a rotating column of fluid in the modes $m=0$, $j=1$ (photographs 1 and 2), $m=0$, $j=2$ (photographs 3, 4, 5, and 6), and $m=0$, $j=3$ (photographs 7 and 8).  The data for the different photographs are as follows: Photographs 1,2: Rotation period 2.295~sec, observed $1/p$ = 1.316, theoretical 1.318, interval between photographs 1.25~sec. Photographs 3,4: 3.947, 0.788, 0.7886, 1.00. Photographs 5,6: 4.040, 0.789, 0.7886, 2.00. Photographs 7,8: 4.058, 0.646, 0.6444, 1.25. (The height of the column is 8.25~cm in all cases.)}

\subsubsection*{(ii) \emph{The case $m = 0$}}
When $m = 0$, the equation governing $\xi_r$ is (cf.\ equation~\eqref{eq:7-58})
\begin{equation}
\label{eq:7-99}
\frac{d^2\xi_r}{dr^2}+\frac{1}{r}\frac{d\xi_r}{dr}
+ \left(a^2\!\left(\frac{4\Omega^2}{p^2}-1\right)
- \frac{1}{r^2}\right)\xi_r = 0.
\end{equation}
The solution of this equation which vanishes at $r = 1$ and $\eta$ can be
written in the form
\begin{equation}
\label{eq:7-100}
\xi_r = \text{constant}\{Y_1(\alpha\eta)J_1(\alpha r) - J_1(\alpha\eta)Y_1(\alpha r)\},
\end{equation}
where $\alpha$ is a root of the equation
\begin{equation}
\label{eq:7-101}
Y_1(\alpha\eta)J_1(\alpha) - J_1(\alpha\eta)Y_1(\alpha) = 0.
\end{equation}
If $\alpha_j$ ($j = 1, 2, \ldots$) denote the roots of this equation, then
\begin{equation}
\label{eq:7-102}
a^2(4\Omega^2/p^2-1) = \alpha_j^2
\end{equation}
or
\begin{equation}
\label{eq:7-103}
p = \pm\frac{2\Omega}{\sqrt{1+\alpha_j^2/a^2}}.
\end{equation}

The values of $\alpha_j$ for a few values of $\eta$ are given in Table~XXXI.  For
$\eta = 0$ and $0.5$, the values of $\alpha_2$ and $\alpha_3$ are also given.

\subsection*{(b) \emph{The case $\Omega = A+B/r^4$ and $m=0$}}
A further case of some interest is when the distribution of the angular
velocity is that which obtains in a viscous liquid.  In that case (cf.\
equation~\eqref{eq:7-20}),
\begin{equation}
\label{eq:7-104}
\Phi = 4A(A+B/r^4);
\end{equation}

\begin{table}[htbp]
\centering
\caption{The roots $\alpha_j$ of equation~\eqref{eq:7-101} for some values of $\eta$}
\label{tab:XXXI}
\begin{tabular}{cccc}
\hline
$\eta$ & $\alpha_1$ & $\alpha_2$ & $\alpha_3$ \\
\hline
$0{\cdot}0$ & $3{\cdot}8317$ & $7{\cdot}01559$ & $10{\cdot}17347$ \\
$0{\cdot}2$ & $4{\cdot}2357$ & & \\
$0{\cdot}3$ & $4{\cdot}79578$ & & \\
$0{\cdot}4$ & $5{\cdot}30116$ & & \\
$0{\cdot}5$ & $6{\cdot}49316$ & $12{\cdot}64470$ & $18{\cdot}88893$ \\
$0{\cdot}8$ & $7{\cdot}93999$ & & \\
$0{\cdot}95$ & $15{\cdot}73587$ & & \\
\hline
\end{tabular}
\end{table}

and equation~\eqref{eq:7-58} governing $\xi_r$ (in case $m = 0$) can be written in the form
\begin{equation}
\label{eq:7-105}
(DD_* - k^2)\xi_r = -4\frac{k^2}{p^2}\,A\!\left(A+\frac{B}{r^4}\right)\xi_r,
\end{equation}
where
\begin{equation}
\label{eq:7-106}
D = \frac{d}{dr} \quad \text{and} \quad
D_* = \frac{d}{dr}+\frac{1}{r},
\end{equation}
and the boundary conditions are
\begin{equation}
\label{eq:7-107}
\xi_r = 0 \quad \text{for } r = R_1 \text{ and } R_2.
\end{equation}

\subsubsection*{(i) \emph{The solution for a narrow gap}}
Reid has recently shown how equation~\eqref{eq:7-105} can be solved explicitly
in terms of known, tabulated functions in the case when
\begin{equation}
\label{eq:7-108}
d = (R_2-R_1) \ll \tfrac{1}{2}(R_2+R_1).
\end{equation}
When~\eqref{eq:7-108} obtains, we need not distinguish between $D$ and $D_*$ and we
can further replace $A+B/r^4$ which occurs on the right-hand side of
equation~\eqref{eq:7-105} by
\begin{equation}
\label{eq:7-109}
\Omega_1\!\left[1-(1-\mu)\frac{r-R_1}{R_2-R_1}\right].
\end{equation}

In rewriting equation~\eqref{eq:7-105} in the framework of the foregoing approximations, we shall find it convenient to measure radial distances from the
surface of the inner cylinder in the unit $d = R_2-R_1$.  Thus, letting
\begin{equation}
\label{eq:7-110}
\zeta = \frac{r-R_1}{d} \quad \text{and} \quad a = k(R_2-R_1),
\end{equation}
we have to solve
\begin{equation}
\label{eq:7-111}
(D^2-a^2)\xi_r = -a^2\lambda[1-(1-\mu)\zeta]\xi_r,
\end{equation}
where $D$ now stands for $d/d\zeta$ and (cf.\ equations~\eqref{eq:7-3})
\begin{equation}
\label{eq:7-112}
\lambda = \frac{4A\Omega_1}{p^2} = \frac{4\Omega_1^2}{p^2}\,\eta^4\,
\frac{1-\mu/\eta^2}{1-\eta^2},
\end{equation}
and the boundary conditions are
\begin{equation}
\label{eq:7-113}
\xi_r = 0 \quad \text{for } \zeta = 0 \text{ and } 1.
\end{equation}

In the case $\mu = 1$, equation~\eqref{eq:7-111} admits of elementary integration
and the proper solutions which satisfy the boundary conditions~\eqref{eq:7-113} are
given by
\begin{equation}
\label{eq:7-114}
\xi_r = \text{constant}\;\sin n\pi\zeta,
\end{equation}
where $n$ is an integer.  The corresponding characteristic equation is
\begin{equation}
\label{eq:7-115}
\frac{1}{\lambda} = \frac{a^2}{a^2+n^2\pi^2},
\end{equation}
or, since $\mu = 1$,
\begin{equation}
\label{eq:7-116}
p = \frac{2\Omega_1}{\sqrt{1+n^2\pi^2/a^2}};
\end{equation}
this is clearly the asymptotic form of~\eqref{eq:7-103} for $\eta \to 1$.

For $\mu < 1$, we can reduce equation~\eqref{eq:7-111} to the standard form
\begin{equation}
\label{eq:7-117}
\frac{d^2\xi_r}{dx^2} = x\xi_r,
\end{equation}
by the substitution
\begin{equation}
\label{eq:7-118}
x = \left[\frac{a}{\lambda(1-\mu)}\right]^{2/3}\!
\{\lambda - \lambda[1-(1-\mu)\zeta]\}.
\end{equation}
The boundary conditions, then, are
\begin{equation}
\label{eq:7-119}
\xi_r = 0 \quad \text{for } x = x_1 \text{ and } x = x_2,
\end{equation}
where
\begin{equation}
\label{eq:7-120}
x_1 = \left[\frac{a}{\lambda(1-\mu)}\right]^{2/3}\!(1-\lambda)
\quad \text{and} \quad
x_2 = \left[\frac{a}{\lambda(1-\mu)}\right]^{2/3}\!(1-\lambda\mu).
\end{equation}
The general solution of equation~\eqref{eq:7-117} can be expressed in terms of
Bessel functions of order $\frac{1}{3}$, or more conveniently in terms of Airy's
functions\footnote{The definitions of these functions are
\[
\mathrm{Ai}(x) = \tfrac{1}{3}x^{1/2}\{I_{-1/3}(y) - I_{1/3}(y)\}
\]
and
\[
\mathrm{Bi}(x) = (x/3)^{1/2}\{I_{-1/3}(y) + I_{1/3}(y)\},
\]
where $y = \frac{2}{3}x^{3/2}$.} $\mathrm{Ai}(x)$ and $\mathrm{Bi}(x)$; thus
\begin{equation}
\label{eq:7-121}
\xi_r = A\,\mathrm{Ai}(x) + B\,\mathrm{Bi}(x),
\end{equation}
where $A$ and $B$ are constants.  The boundary conditions~\eqref{eq:7-119} lead to
the characteristic equation
\begin{equation}
\label{eq:7-122}
\frac{\mathrm{Ai}(x_1)}{\mathrm{Bi}(x_1)} = \frac{\mathrm{Ai}(x_2)}{\mathrm{Bi}(x_2)}.
\end{equation}
Equation~\eqref{eq:7-122} determines $x_2$ as a function of $x_1$ for the different modes;
the relationship for the lowest mode is illustrated in Fig.~65.

\figurePlaceholder{65}{The $(x_2,x_1)$-relation for the lowest mode
as determined by equation~\eqref{eq:7-122}.}

$(x_2,x_1)$-relationship determined in this fashion, the required relation
between $a$ and $\sqrt{\lambda}$ is given, parametrically, by
\begin{equation}
\label{eq:7-123}
a = (x_2-x_1)\sqrt{\frac{x_2-\mu x_1}{1-\mu}}
\quad \text{and} \quad
\frac{1}{\sqrt{\lambda}} = \sqrt{\frac{x_2-\mu x_1}{x_2-x_1}}.
\end{equation}
The dispersion relations for $\mu = 1$, 0, and $-1$ are illustrated in Fig.~66.
In Fig.~67\,$a$ the limiting form of velocity profile for the lowest mode of
instability when $\mu \to -\infty$ is illustrated; and the corresponding cell
pattern is shown in Fig.~67\,$b$.

\figurePlaceholder{66}{The dispersion relations in the limit of zero viscosity
for $\mu=1$, 0, and $-1$ for the lowest modes of the flow between two cylinders with a narrow gap $d$.  The abscissa gives
the wave number in the unit $1/d$ and the ordinate gives the
growth rate (or the circular frequency) in the unit $\sqrt{4A\Omega_1}$.}

\figurePlaceholder{67}{The limiting form of the velocity profile (a) and the cell pattern (b) for
lowest inviscid mode of instability as $\mu\to-\infty$ and $a/(1-\mu) = 2{\cdot}03$.}

\subsubsection*{(ii) \emph{The formal solution for a wide gap}}
When we wish to allow for the finite difference between $R_1$ and $R_2$,
we must consider equation~\eqref{eq:7-105} without making any approximations.

Now measuring $r$ in units of $R_2$ and letting $a = kR_2$, we have the
equation (cf.\ equations~\eqref{eq:7-22} and~\eqref{eq:7-23})
\begin{equation}
\label{eq:7-124}
(DD_*-a^2)\xi_r = a^2\frac{4\Omega_1^2}{p^2}\,\eta^4\,
\frac{(1-\mu)(1-\mu/\eta^2)}{(1-\eta^2)^2}
\!\left(\frac{1}{r^4}-\kappa\right)\xi_r.
\end{equation}

Letting
\begin{equation}
\label{eq:7-125}
Q_1^2 = -\frac{4\Omega_1^2}{p^2}\,\eta^4\,
\frac{(1-\mu)(1-\mu/\eta^2)}{(1-\eta^2)^2}
\end{equation}
and
\begin{equation}
\label{eq:7-126}
Q_1^2\kappa = -Q_1^2\,\kappa,
\end{equation}
we can rewrite equation~\eqref{eq:7-124} in the form
\begin{equation}
\label{eq:7-127}
\frac{d^2\xi_r}{dr^2}+\frac{1}{r}\frac{d\xi_r}{dr}
+ \left(a^2 Q_1^2 - 1\right)\frac{1-a^2 Q_1^2}{r^4}\,\xi_r = 0.
\end{equation}
The general solution of this equation is
\begin{equation}
\label{eq:7-128}
\xi_r = \mathscr{C}_\nu(ar),
\end{equation}
where
\begin{equation}
\label{eq:7-129}
\nu = \sqrt{1-a^2 Q_1^2}
\quad \text{and} \quad
\alpha = a\sqrt{Q_1^2-1},
\end{equation}
and $\mathscr{C}_\nu$ represents a general cylinder function of order $\nu$.

The boundary conditions require $\mathscr{C}_\nu(\alpha r)$ to vanish at $r = 1$ and $\eta$;
and these conditions will determine $\alpha$.  The required solution can, in fact,
be expressed in the form
\begin{equation}
\label{eq:7-130}
\xi_r = \text{constant}\{J_{-\nu}(\alpha\eta)J_\nu(\alpha r)
- J_\nu(\alpha\eta)J_{-\nu}(\alpha r)\}.
\end{equation}
Defined in this manner, $\xi_r$ clearly vanishes at $r = \eta$; and the condition
that it also vanishes at $r = 1$ leads to the equation
\begin{equation}
\label{eq:7-131}
J_{-\nu}(\alpha\eta)J_\nu(\alpha)
- J_\nu(\alpha\eta)J_{-\nu}(\alpha) = 0.
\end{equation}

With $\alpha$ determined as a root of this equation for some assigned $\nu$, the
required relation between $a^2$ and $p^2$ follows from equations~\eqref{eq:7-126} and~\eqref{eq:7-129}.

From the general theory of \S\S~66 and 67, we know that
\begin{equation}
\label{eq:7-132}
p^2 > 0 \text{ for } \mu > \eta^2 \quad \text{and} \quad
p^2 < 0 \text{ for } \mu < \eta^2.
\end{equation}
Therefore,
\begin{equation}
\label{eq:7-133}
Q_1^2 > 0 \quad \text{or} \quad < 0 \quad \text{according as} \quad
\mu < 1 \quad \text{or} \quad > 1.
\end{equation}
But $Q_1^2$ is positive or negative according as $p^2$ is positive or negative,
i.e.\ according as $\mu > \eta^2$ or $\mu < \eta^2$.  Thus, for the unstable modes, $Q_1^2 < 0$
and $\alpha$ becomes imaginary.  If $\nu$ were real, $\mathscr{C}_\nu(\alpha r)$ would become a linear
combination of $I_\nu(|\alpha|r)$ and $K_\nu(|\alpha|r)$.  But no linear combination of
$I_\nu(x)$ and $K_\nu(x)$ can vanish at two specified points.  Therefore, for the
unstable modes, $\nu$ must be imaginary and the cylinder functions, in terms
of which the solutions are expressed, are of imaginary order and argument.

\section{On viscous Couette flow}\label{sec:7-69}
We now turn our attention to stationary viscous flow between two
rotating coaxial cylinders.

In cylindrical polar coordinates, the Navier--Stokes equations for
viscous incompressible fluids take the forms:
\begin{equation}
\label{eq:7-134}
\frac{\partial u_r}{\partial t}
+ (\mathbf{u}.\mathrm{grad})\,u_r - \frac{u_\theta^2}{r}
= -\frac{\partial}{\partial r}\!\left(\frac{p}{\rho}\right)
+ \nu\!\left(\nabla^2 u_r - \frac{2}{r^2}\frac{\partial u_\theta}{\partial \theta}
- \frac{u_r}{r^2}\right),
\end{equation}
\begin{equation}
\label{eq:7-135}
\frac{\partial u_\theta}{\partial t}
+ (\mathbf{u}.\mathrm{grad})\,u_\theta + \frac{u_r u_\theta}{r}
= -\frac{1}{r}\frac{\partial}{\partial r}\!\left(\frac{p}{\rho}\right)
+ \nu\!\left(\nabla^2 u_\theta + \frac{2}{r^2}\frac{\partial u_r}{\partial \theta}
- \frac{u_\theta}{r^2}\right),
\end{equation}
\begin{equation}
\label{eq:7-136}
\frac{\partial u_z}{\partial t}
+ (\mathbf{u}.\mathrm{grad})\,u_z
= -\frac{\partial}{\partial z}\!\left(\frac{p}{\rho}\right)
+ \nu \nabla^2 u_z,
\end{equation}
where
\begin{equation}
\label{eq:7-137}
\mathbf{u}.\mathrm{grad} = u_r \frac{\partial}{\partial r}
+ \frac{u_\theta}{r}\frac{\partial}{\partial \theta}
+ u_z \frac{\partial}{\partial z}
\end{equation}
and
\begin{equation}
\label{eq:7-138}
\nabla^2 = \frac{\partial^2}{\partial r^2}
+ \frac{1}{r}\frac{\partial}{\partial r}
+ \frac{1}{r^2}\frac{\partial^2}{\partial \theta^2}
+ \frac{\partial^2}{\partial z^2}.
\end{equation}
We have also the equation of continuity,
\begin{equation}
\label{eq:7-139}
\frac{\partial u_r}{\partial r}
+ \frac{u_r}{r}
+ \frac{1}{r}\frac{\partial u_\theta}{\partial \theta}
+ \frac{\partial u_z}{\partial z} = 0.
\end{equation}
These equations allow a stationary solution of the form
\begin{equation}
\label{eq:7-140}
u_r = u_z = 0 \quad \text{and} \quad u_\theta = V(r),
\end{equation}
provided
\begin{equation}
\label{eq:7-141}
\frac{d}{dr}\!\left(\frac{p}{\rho}\right) = \frac{V^2}{r}
\end{equation}
and
\begin{equation}
\label{eq:7-142}
\nu\!\left(\nabla^2 V - \frac{V}{r^2}\right)
= \nu\,\frac{d}{dr}\!\left(\frac{d}{dr}+\frac{1}{r}\right)V = 0.
\end{equation}
From equation~\eqref{eq:7-142}, it is apparent that the most general form of $V(r)$
which is compatible with $\nu \neq 0$ is
\begin{equation}
\label{eq:7-143}
V = Ar + \frac{B}{r},
\end{equation}
where $A$ and $B$ are two arbitrary constants.  The corresponding expression for the angular velocity is
\begin{equation}
\label{eq:7-144}
\Omega = A + \frac{B}{r^2}.
\end{equation}
This is the solution which was quoted in \S~65.  And as we remarked at
that time, the appearance of two arbitrary constants in the solution~\eqref{eq:7-144}
corresponds to the possibility of assigning to the two cylinders, confining
the fluid, angular velocities at will.  The constants $A$ and $B$ can, therefore, be related to the angular velocities $\Omega_1$ and $\Omega_2$ with which the two
cylinders are rotated.  We have (cf.\ equations~\eqref{eq:7-3} and~\eqref{eq:7-4})
\begin{equation}
\label{eq:7-145}
A = -\Omega_1\,\eta^2\,\frac{1-\mu/\eta^2}{1-\eta^2}
\quad \text{and} \quad
B = \Omega_1\,\frac{R_1^2(1-\mu)}{1-\eta^2},
\end{equation}
where
\begin{equation}
\label{eq:7-146}
\mu = \Omega_2/\Omega_1 \quad \text{and} \quad \eta = R_1/R_2.
\end{equation}
As we have seen in \S~66, Rayleigh's criterion for the stability of inviscid
Couette flow applied to the distribution~\eqref{eq:7-144} yields
\begin{equation}
\label{eq:7-147}
\mu > \eta^2 \quad \text{for stability.}
\end{equation}
On general grounds, we should expect the effect of viscosity is to
postpone the onset of instability; and that the new criterion for stability
must essentially take the form of specifying the extent to which the
condition~\eqref{eq:7-147} can be violated before instability will manifest itself.

The precise form which the criterion takes was first investigated, both
theoretically and experimentally, by G.~I.~Taylor.  In the case when the
gap $R_2-R_1$ between the two cylinders can be considered as small
compared to the mean radius $\frac{1}{2}(R_1+R_2)$, Taylor found an explicit
analytical expression for the criterion; and he was able to confirm by
experiments that the marginal state is stationary and exhibits a break-up
of the basic flow into a cellular pattern.  Fig.~68, which is a reproduction of one of Taylor's photographs, clearly shows the manner of the
onset of instability.  In \S~74 we return to a more detailed account of the
experimental work which has been carried out since Taylor's early
pioneering work.

\figurePlaceholder{68}{One of G.~I.~Taylor's photographs illustrating the onset of instability
in the flow between rotating cylinders.  The radii of the two cylinders in this
case are 4.035~cm and 3.25~cm respectively; the cylinders are rotating in the
same direction.}

\section{The perturbation equations}\label{sec:7-70}
We shall now investigate the stability of the flow described by equations~\eqref{eq:7-141} and~\eqref{eq:7-143}.  Let the perturbed state be characterized by
\begin{equation}
\label{eq:7-148}
u_r, \quad V\!+\!u_\theta, \quad u_z, \quad \text{and} \quad
\delta p/\rho = \varpi.
\end{equation}
Assuming that the various perturbations are axisymmetric\footnote{The general non-axisymmetric case has not been investigated for this problem.} and independent of $\theta$, we obtain from~\eqref{eq:7-134}--\eqref{eq:7-136} the linearized equations
\begin{equation}
\label{eq:7-149}
\frac{\partial u_r}{\partial t}
- \frac{2V}{r}\, u_\theta
= -\frac{\partial \varpi}{\partial r}
+ \nu\!\left(\nabla^2 u_r - \frac{u_r}{r^2}\right),
\end{equation}
\begin{equation}
\label{eq:7-150}
\frac{\partial u_\theta}{\partial t}
+ \left(\frac{dV}{dr}+\frac{V}{r}\right) u_r
= \nu\!\left(\nabla^2 u_\theta - \frac{u_\theta}{r^2}\right),
\end{equation}
and
\begin{equation}
\label{eq:7-151}
\frac{\partial u_z}{\partial t}
= -\frac{\partial \varpi}{\partial z}
+ \nu \nabla^2 u_z,
\end{equation}
where $\nabla^2$ has now the meaning
\begin{equation}
\label{eq:7-152}
\nabla^2 = \frac{\partial^2}{\partial r^2}
+ \frac{1}{r}\frac{\partial}{\partial r}
+ \frac{\partial^2}{\partial z^2}.
\end{equation}
Also, for axisymmetric motions, the equation of continuity reduces to
\begin{equation}
\label{eq:7-153}
\frac{\partial u_r}{\partial r}
+ \frac{u_r}{r}
+ \frac{\partial u_z}{\partial z} = 0.
\end{equation}

By analysing the disturbance into normal modes, we seek solutions
of the foregoing equations which are of the forms
\begin{equation}
\label{eq:7-154}
\left.\begin{aligned}
u_r &= e^{pt}u(r)\cos kz;\quad &
u_\theta &= e^{pt}v(r)\sin kz\\
u_z &= e^{pt}w(r)\cos kz;\quad &
\varpi &= e^{pt}\varpi(r)\cos kz
\end{aligned}\right\},
\end{equation}
where $k$ is the wave number of the disturbance in the axial direction
and $p$ is a constant which can be complex.  For solutions of the form~\eqref{eq:7-154}, equations~\eqref{eq:7-149}--\eqref{eq:7-153} become

\begin{equation}
\label{eq:7-155}
\nu\!\left(DD_*-k^2-\frac{p}{\nu}\right)u
+ 2\frac{V}{r}\,v = \frac{d\varpi}{dr},
\end{equation}
\begin{equation}
\label{eq:7-156}
\nu\!\left(D_* D - k^2 - \frac{p}{\nu}\right)v - (D_*V)\,u = 0,
\end{equation}
\begin{equation}
\label{eq:7-157}
\nu\!\left(D_* D-k^2-\frac{p}{\nu}\right)w = -k\varpi,
\end{equation}
\begin{equation}
\label{eq:7-158}
\nabla^2 = \left(\frac{d}{dr}+\frac{1}{r}\right)\frac{1}{r}\frac{d}{dr}\, r - k^2
= D_* D - k^2,
\end{equation}
and
\begin{equation}
\label{eq:7-159}
D_* u = -kw.
\end{equation}
[In the foregoing equations, $D$ and $D_*$ have the same meanings as in
\S~68\,($b$), equation~\eqref{eq:7-106}.]

Eliminating $w$ between equations~\eqref{eq:7-157} and~\eqref{eq:7-159}, we have
\begin{equation}
\label{eq:7-160}
\frac{\nu}{k^2}\!\left(D_* D-k^2-\frac{p}{\nu}\right)D_*\,u = \varpi.
\end{equation}
Inserting this expression for $\varpi$ in equation~\eqref{eq:7-155}, we find after some
rearranging,
\begin{equation}
\label{eq:7-161}
\frac{\nu}{k^2}\!\left(DD_*-k^2-\frac{p}{\nu}\right)(DD_*-k^2)u
= 2\frac{V}{r}\,v.
\end{equation}
This equation must be considered together with
\begin{equation}
\label{eq:7-162}
\nu\!\left(DD_*-k^2-\frac{p}{\nu}\right)v = (D_*V)u.
\end{equation}
These equations are general and do not depend on any particular form
of $V(r)$.

Measuring $r$ in units of the radius $R_2$ of the outer cylinder and writing
\begin{equation}
\label{eq:7-163}
k^2 = a^2/R_2^2 \quad \text{and} \quad
\sigma = p R_2^2/\nu,
\end{equation}
equations~\eqref{eq:7-161} and~\eqref{eq:7-162} become (when $V(r)$ has the particular form~\eqref{eq:7-143}):
\begin{equation}
\label{eq:7-164}
(DD_*-a^2)(DD_*-a^2)u = a^2\frac{2B}{\nu}\!\left(\frac{1}{r^2}+\frac{A\,R_2^4}{B}\right)v,
\end{equation}
and
\begin{equation}
\label{eq:7-165}
(DD_*-a^2-\sigma)v = \frac{2A}{\nu}\,R_2^2\, u.
\end{equation}
It is convenient to make the transformation
\begin{equation}
\label{eq:7-166}
\frac{2A\,R_2^2}{\nu}\,u \to u;
\end{equation}
the equations then take the more convenient forms
\begin{equation}
\label{eq:7-167}
(DD_*-a^2-\sigma)(DD_*-a^2)u = -Ta^2\!\left(\frac{1}{r^2}-\kappa\right)v,
\end{equation}
and
\begin{equation}
\label{eq:7-168}
(DD_*-a^2-\sigma)v = u,
\end{equation}
where (cf.\ equations~\eqref{eq:7-21}, \eqref{eq:7-23}, and~\eqref{eq:7-32})
\begin{equation}
\label{eq:7-169}
T = -\frac{4AB}{R_2^4} = \frac{4\Omega_1^2\,R_2^4}{\nu^2}\,
\frac{(1-\mu)(1-\mu/\eta^2)}{(1-\eta^2)^2}
\end{equation}
and
\begin{equation}
\label{eq:7-170}
\kappa = -\frac{A\,R_2^4}{B} = \frac{1-\mu/\eta^2}{1-\mu}.
\end{equation}

Solutions of equations~\eqref{eq:7-167} and~\eqref{eq:7-168} must be sought which satisfy
the boundary conditions appropriate for no slip on the cylindrical walls
at $r = 1$ and $\eta$.  These conditions are that all three components of the
velocity vanish on the walls; thus,
\begin{equation}
\label{eq:7-171}
u = v = 0 \quad \text{and} \quad Du = 0
\quad \text{for } r = 1 \text{ and } \eta,
\end{equation}
where the last of the three boundary conditions is equivalent to $w = 0$
(cf.\ equation~\eqref{eq:7-159}).

\subsection*{(a) \emph{The stability of the flow for $\mu > \eta^2$}}
We shall now show that when Rayleigh's criterion $\mu > \eta^2$ is satisfied,
the flow is indeed stable.

First, we may notice certain elementary integral properties of the
operator $DD_*$.  If $f(r)$ and $g(r)$ are any two functions and if one of them,
say $f(r)$, vanishes at the limits of integration,
\begin{equation}
\label{eq:7-172}
\int f\, DD_*\, g\, dr
= -\int\!\left(r\frac{df}{dr}\frac{dg}{dr}+\frac{fg}{r}\right)dr,
\end{equation}
and if the derivative of $f$ also vanishes at the limits,
\begin{equation}
\label{eq:7-173}
\int f\, DD_*\, g\, dr = \int \eta\, DD_*\, f\, dr.
\end{equation}
These relations follow by successive integrations by parts.  Thus, by
writing
\begin{equation}
\label{eq:7-174}
\int f\, DD_*\, g\, dr
= \int\!\left\{
\frac{1}{r}\frac{d}{dr}\!\left(r\frac{dg}{dr}\right)
- \frac{g}{r}\right\} f\, dr,
\end{equation}
the truth of~\eqref{eq:7-172} becomes self-evident; and a further integration by parts
leads to~\eqref{eq:7-173}.

Now returning to equations~\eqref{eq:7-167} and~\eqref{eq:7-168}, multiply~\eqref{eq:7-167} by $ru^*$
and integrate over the range of $r$.  We have
\begin{equation}
\label{eq:7-175}
\int ru^*\{(DD_*-a^2)^2u - \sigma(DD_*-a^2)u\}\,dr
= -\mathscr{T}a^2 \int r\phi(r)vu^*\,dr,
\end{equation}
where, for brevity, we have written
\begin{equation}
\label{eq:7-176}
\mathscr{T} = \frac{T}{a^2} = \frac{4\Omega_1^2\,R_2^4\,(1-\mu/\eta^2)}{\nu^2\,(1-\eta^2)^2}
\quad \text{and} \quad
\phi = (1-\mu)\!\left(\frac{1}{r^2}-\kappa\right).
\end{equation}
Since $u$ and its derivative vanish at $r = 1$ and $\eta$, the integrals on the
left-hand side of~\eqref{eq:7-175} can be transformed to positive definite forms by
making use of~\eqref{eq:7-172} and~\eqref{eq:7-173}.  Thus,
\begin{equation}
\label{eq:7-177}
\int ru^*\{(DD_*-a^2)^2u - \sigma(DD_*-a^2)u\}\,dr
= \int r\{(DD_*-a^2)u\}^2\,dr
+ \left(\frac{1}{r}+a^2 r\right)|u|^2\,dr.
\end{equation}
(This is not quite right as written. The full expansion gives separate integral terms.)

Next, substituting for $u^*$ (from equation~\eqref{eq:7-168}) in the integrand on
the right-hand side of equation~\eqref{eq:7-175}, we obtain
\begin{equation}
\label{eq:7-178}
\int r\phi(r)vu^*\,dr
= \int r\phi(r) v(DD_*-a^2-\sigma^*)v^*\,dr
+ \int r\phi(r)vDD_*v^*\,dr.
\end{equation}
Again, by making use of~\eqref{eq:7-172}, we have
\begin{equation}
\label{eq:7-179}
\int r\phi(r)v DD_*v^*\,dr
= -\int \phi(r)\!\left(r\left|\frac{dv}{dr}\right|^2
+ \frac{|v|^2}{r}\right)dr
+ 2(1-\mu)\int \frac{v}{r}\frac{dv^*}{dr}\,dr.
\end{equation}
Now combining equations~\eqref{eq:7-175}, \eqref{eq:7-177}, \eqref{eq:7-178}, and~\eqref{eq:7-179}, we obtain
\begin{equation}
\label{eq:7-180}
\sigma I_1 + I_2 = \mathscr{T}a^2(a^2+\sigma^*)I_3 + I_4\},
\end{equation}
where
\begin{equation}
\label{eq:7-181}
I_1 = \int \phi(r)\!\left(r\left|\frac{du}{dr}\right|^2
+ \left(\frac{1}{r}+a^2 r\right)|u|^2\right)dr,
\end{equation}
\begin{equation}
\label{eq:7-182}
I_2 = \int |(DD_*-a^2)u|^2\,r\,dr,
\end{equation}
\begin{equation}
\label{eq:7-183}
I_3 = \int \phi(r)|v|^2\,dr,
\end{equation}
and
\begin{equation}
\label{eq:7-184}
I_4 = \int \phi(r)\!\left(r\left|\frac{dv}{dr}\right|^2
+ \frac{|v|^2}{r}\right)dr
- 2(1-\mu)\int \frac{v}{r}\frac{dv^*}{dr}\,dr.
\end{equation}
The integrals $I_2$ and $I_3$ are clearly positive definite.  For $\mu > 0$, $\phi(r) > 0$
(see equation~\eqref{eq:7-26}); so that in this case, $I_1$ is also positive definite.  The
first of the two integrals included in $I_4$ is positive definite for $\mu > 0$; but
the second is complex.  However, the real part of $I_4$ is positive definite
for $\mu > 0$; in fact,
\begin{equation}
\label{eq:7-185}
\mathrm{re}(I_4) = \int r\phi(r)\left|\frac{dv}{dr}
- \frac{v}{r}\right|^2 dr.
\end{equation}
For, expanding the integrand in~\eqref{eq:7-185}, we have
\begin{equation}
\label{eq:7-186}
\int r\phi(r)\left|\frac{dv}{dr}
- \frac{v}{r}\right|^2 dr
= \int \phi(r)\!\left(r\left|\frac{dv}{dr}\right|^2
+ \frac{|v|^2}{r}\right)dr
- \int \phi(r)\frac{d}{dr}|v|^2\,dr;
\end{equation}
but
\begin{equation}
\label{eq:7-187}
\int \phi(r)\frac{d}{dr}|v|^2\,dr
= (1-\mu)\int \frac{1}{r}\!\left(\frac{1}{r^2}-\kappa\right)|v|^2\,dr
= (1-\mu)\int \frac{1}{r^2}\frac{d|v|^2}{dr}\,dr.
\end{equation}
Therefore, the right-hand side of~\eqref{eq:7-186} is, indeed, the real part of $I_4$.
Returning to equation~\eqref{eq:7-180} and equating the real parts of this
equation, we obtain
\begin{equation}
\label{eq:7-188}
\mathrm{re}(\sigma)\{I_1-\mathscr{T}a^2 I_3+I_4-\mathscr{T}a^2 \sigma^* I_3 + \mathrm{re}(I_4)\} = 0.
\end{equation}
When $\mu > \eta^2$, $\mathscr{T} < 0$ and the coefficient of $\mathrm{re}(\sigma)$ in equation~\eqref{eq:7-188} is
positive definite; and so also are the remaining terms in the equation.
Therefore,
\begin{equation}
\label{eq:7-189}
\mathrm{re}(\sigma) < 0 \quad \text{for } \mu > \eta^2.
\end{equation}
and the flow is stable; this result is entirely to be expected on physical
grounds.  Nevertheless, it appears to be the only one which can be
established by general analytical arguments.  In particular, it does not
seem that one can deduce the general validity of the principle of the
exchange of stabilities for this problem.  For example, by equating the
imaginary parts of equation~\eqref{eq:7-180}, we obtain (cf.\ equation~\eqref{eq:7-176})
\begin{equation}
\label{eq:7-190}
\mathrm{im}(\sigma)(I_1+\mathscr{T}a^2 I_3)
= -2\mathscr{T}a^2\,\mathrm{im}\int\frac{v}{r}\frac{dv^*}{dr}\,dr,
\end{equation}
and no general conclusions can be drawn from this equation; when $\mu < 0$,
even $I_4$ is not positive definite!

\section{The solution for the case of a narrow gap when the marginal
state is stationary}\label{sec:7-71}
If the gap $R_2-R_1$ between the two cylinders is small compared to
their mean radius $\frac{1}{2}(R_1+R_2)$, we need not (as in \S~68\,($b$)) distinguish
between $D$ and $D_*$ in equations~\eqref{eq:7-161} and~\eqref{eq:7-162}; and we can also replace
$(A+B/r^4)$ which occurs on the right-hand side of equation~\eqref{eq:7-161} by
\begin{equation}
\label{eq:7-191}
\Omega_1\!\left[1-(1-\mu)\frac{r-R_1}{R_2-R_1}\right].
\end{equation}

In rewriting equations~\eqref{eq:7-161} and~\eqref{eq:7-162} in the framework of these
approximations, it will be convenient to measure radial distances from the
surface of the inner cylinder in the unit $d = R_2-R_1$.  Thus, letting
\begin{equation}
\label{eq:7-192}
\zeta = (r-R_1)/d, \quad k = a/d, \quad \text{and} \quad
\sigma = pd^2/\nu,
\end{equation}
we have to consider the equations
\begin{equation}
\label{eq:7-193}
(D^2-a^2-\sigma)(D^2-a^2)u = -Ta^2[1-(1-\mu)\zeta]v
\end{equation}
and
\begin{equation}
\label{eq:7-194}
(D^2-a^2-\sigma)v = \frac{2Ad^2}{\nu}\,u.
\end{equation}
By the further transformation
\begin{equation}
\label{eq:7-195}
u \to \frac{2\Omega_1\,d^2 a^2}{\nu}\,u,
\end{equation}
the equations become
\begin{equation}
\label{eq:7-196}
(D^2-a^2-\sigma)(D^2-a^2)u = (1+\alpha\zeta)v
\end{equation}
and
\begin{equation}
\label{eq:7-197}
(D^2-a^2-\sigma)v = -Ta^2 u,
\end{equation}
where, now,
\begin{equation}
\label{eq:7-198}
T = \frac{4A\Omega_1}{\nu^2}\,d^4
\end{equation}
and
\begin{equation}
\label{eq:7-199}
\alpha = -(1-\mu).
\end{equation}
Equations~\eqref{eq:7-196} and~\eqref{eq:7-197} must be considered together with the boundary
conditions
\begin{equation}
\label{eq:7-200}
u = Du = v = 0 \quad \text{for } \zeta = 0 \text{ and } 1.
\end{equation}
We are primarily interested in the solutions of equations~\eqref{eq:7-196} and~\eqref{eq:7-197} (subject to the boundary conditions~\eqref{eq:7-200}) for various values of $\alpha$
for which the real part of $\sigma$ is zero.  The method described in Chapter~III,
\S~31 in a different connexion is applicable to this problem.  Thus, we
must obtain solutions for two cases: when $\sigma$ is zero and the marginal
state is stationary and when $\sigma$ is imaginary and the marginal state is
oscillatory.  In the latter case for each value of $\alpha$, $\sigma$ must be determined
by the condition that $T$ is real.\footnote{It is, of course, possible that under certain circumstances solutions with this property
do not exist.}  In either case, we must find the minimum of $T$ as a function of $a$; and depending on which of the two minima
is lower, we shall have the onset of instability as a stationary secondary
flow or as overstability.  Careful experiments on the onset of instability
by Taylor and others have failed to reveal any suggestions of overstability.  For this reason, the case $\sigma = 0$ is the only one which has been
considered in the literature.  However, as no general arguments for the
validity of the principle of the exchange of stabilities have been found
for this problem, the case of overstability requires investigation.  We
return to this question in \S~72.

When the marginal state is stationary, the equations to be solved are
\begin{equation}
\label{eq:7-201}
(D^2-a^2)^2 u = (1+\alpha\zeta)v
\end{equation}
and
\begin{equation}
\label{eq:7-202}
(D^2-a^2)v = -Ta^2 u,
\end{equation}
together with the boundary conditions~\eqref{eq:7-200}.

\subsection*{(a) \emph{The solution of the characteristic value problem for the case
$\alpha = 0$}}
It can be readily verified that the characteristic value problem presented by equations~\eqref{eq:7-201}--\eqref{eq:7-202} is not self-adjoint in the usual sense.
For this reason, the method to be described below was patterned after the
ones which have been found successful in cases where the problems are
self-adjoint.  However, Roberts has recently found a variational basis
for the method; this is considered in Appendix~IV.

The method of solution we shall adopt is the following.

Since $v$ is required to vanish at $\zeta = 0$ and $1$, we expand it in a sine
series of the form
\begin{equation}
\label{eq:7-203}
v = \sum_{m=1}^{\infty} C_m \sin m\pi\zeta.
\end{equation}
Having chosen $v$ in this manner, we next \emph{solve} the equation,
\begin{equation}
\label{eq:7-204}
(D^2-a^2)^2 u = (1+\alpha\zeta)\sum_{m=1}^{\infty} C_m \sin m\pi\zeta,
\end{equation}
obtained by inserting~\eqref{eq:7-203} in~\eqref{eq:7-201}, and arrange that the solution
satisfies the four remaining boundary conditions on $u$.  With $u$ determined
in this fashion and $v$ given by~\eqref{eq:7-203}, equation~\eqref{eq:7-202} will lead, as we shall
presently see, to a secular equation for $T$.

The solution of equation~\eqref{eq:7-204} is straightforward.  The general
solution can be written in the form
\begin{multline}
\label{eq:7-205}
u = \sum_{m=1}^{\infty} \frac{C_m}{m^2\pi^2+a^2}
\biggl\{
A_1^{(m)}\cosh a\zeta + A_2^{(m)}\,\zeta\cosh a\zeta
+ A_1^{(m)}\,\zeta^2\cosh a\zeta + \\
+ B_1^{(m)}\,\zeta\sinh a\zeta
+ A_2^{(m)}\,\zeta\sinh a\zeta
+ B_2^{(m)}\,\zeta^2\sinh a\zeta + \\
+ (1+\alpha\zeta)\sin m\pi\zeta
+ \frac{4\alpha m\pi}{m^2\pi^2+a^2}\cos m\pi\zeta
\biggr\},
\end{multline}
where $A_1^{(m)}, \ldots, B_2^{(m)}$ are constants of integration to be
determined by the boundary conditions~\eqref{eq:7-218}.  These latter conditions lead to the equations:
\begin{align}
\label{eq:7-206}
A_1^{(m)} &= -\frac{4m\pi\alpha}{m^2\pi^2+a^2}, \nonumber\\[4pt]
B_1^{(m)} &= \frac{m\pi}{m^2\pi^2+a^2}
\{\alpha + \beta_m(\sinh a + a\cosh a) - \gamma_m\sinh a\}, \nonumber\\[4pt]
A_2^{(m)} &= -\frac{m\pi}{m^2\pi^2+a^2}
\{(\sinh^2 a + a)(\sinh a + a\cosh a) - \gamma_m \sinh a\}, \nonumber\\[4pt]
B_2^{(m)} &= \frac{m\pi}{m^2\pi^2+a^2}
\{(\sinh a\cosh a - a)(1 + \alpha)(-1)^{m+1} - \\
&\qquad -\gamma_m(\alpha\cosh a - \alpha\cosh a)\},
\end{align}
where
\begin{equation}
\label{eq:7-207}
\Delta = \sinh^2 a - a^2,
\end{equation}
\begin{equation}
\beta_m = \frac{4\alpha}{m^2\pi^2+a^2}\{(-1)^{m+1}\cosh a\},
\end{equation}
and
\begin{equation}
\gamma_m = (-1)^{m+1}(1+\alpha) + \frac{4\alpha}{m^2\pi^2+a^2}\, a\sinh a.
\end{equation}

Now substituting for $v$ and $u$ from equations~\eqref{eq:7-203} and~\eqref{eq:7-205} in equation~\eqref{eq:7-202}, we obtain
\begin{multline}
\label{eq:7-209}
\sum_{n=1}^{\infty} C_n(n^2\pi^2+a^2)\sin n\pi\zeta \\
= Ta^2 \sum_{m=1}^{\infty}\frac{C_m}{m^2\pi^2+a^2}
\biggl\{
A_1^{(m)}\cosh a\zeta + B_1^{(m)}\sinh a\zeta
+ A_2^{(m)}\,\zeta\cosh a\zeta + \\
+ B_2^{(m)}\,\zeta\sinh a\zeta
+ (1+\alpha\zeta)\sin m\pi\zeta
+ \frac{4\alpha m\pi}{m^2\pi^2+a^2}\cos m\pi\zeta
\biggr\}.
\end{multline}

Multiplying equation~\eqref{eq:7-209} by $\sin n\pi\zeta$ and integrating over the range of
$\zeta$, we obtain a system of linear homogeneous equations for the constants
$\mathscr{C}_m = C_m(m^2\pi^2+a^2)^{1/2}$; and the requirement that these constants are
not all zero leads to the secular equation
\begin{multline}
\label{eq:7-210}
\left|\frac{n\pi}{(n^2\pi^2+a^2)^{1/2}}
\{[1+(-1)^{n+1}\cosh a]\,A_1^{(m)}+\right. \\
+(-1)^{n+1}\left[\cosh a
- \frac{2a}{n^2\pi^2+a^2}\,\sinh a\right]\,A_2^{(m)} + \\
+\left[(-1)^{n+1}\sinh a
- \frac{2a}{n^2\pi^2+a^2}
\{1+(-1)^{n+1}\cosh a\}\right] B_2^{(m)} + \\
\left.+\alpha X_{mn}+\tfrac{1}{2}\delta_{mn}
- \tfrac{1}{2}(n^2\pi^2+a^2)\frac{\delta_{mn}}{a^2}\right| = 0,
\end{multline}
where
\begin{equation}
\label{eq:7-211}
X_{mn} = \begin{cases}
0 & \text{if } m+n \text{ is even and } m \neq n,\\
\frac{1}{2} & \text{if } m = n,\\
\dfrac{4mn}{n^2-m^2}\left(\dfrac{3}{m^2\pi^2+a^2}
- \dfrac{2}{n^2-m^2}\right.\cdot
\left.\dfrac{1}{\pi(n^2-m^2)}\right)
& \text{if } m+n \text{ is odd.}
\end{cases}
\end{equation}

On using the first two equations of~\eqref{eq:7-206}, equation~\eqref{eq:7-210} simplifies to
the form
\begin{multline}
\label{eq:7-212}
\left|\frac{n\pi}{n^2\pi^2+a^2}\frac{4mm\pi x}{m^2\pi^2+a^2}
\{(-1)^{m+n}-1\} - \right.\\
-\frac{2a}{n^2\pi^2+a^2}
\{(-1)^{n+1}\{A_2^{(m)}\sinh a + B_1^{(m)}\cosh a\} +
B_2^{(m)}\} + \\
\left.+\alpha X_{mn}+\tfrac{1}{2}\delta_{mn}
-\tfrac{1}{2}(n^2\pi^2+a^2)\frac{\delta_{mn}}{a^2}\right| = 0;
\end{multline}
and on substituting for the constants $A_2^{(m)}$ and $B_2^{(m)}$ their explicit solutions
given in~\eqref{eq:7-207}, we find that equation~\eqref{eq:7-212} simplifies greatly and we are
left with
\begin{multline}
\label{eq:7-213}
\left|\frac{4mm'n^2\alpha}{(n^2\pi^2+a^2)(m^2\pi^2+a^2)}
\{(-1)^{m+n}-1\} -\right.\\
- \frac{2amm'\pi^2\alpha}{(m^2\pi^2+a^2)^2(\sinh a)^2}
\{(\sinh a\cosh a - a)[1+(1+\alpha)(-1)^{m+n}] + \\
+(\sinh a - a\cosh a)[(-1)^{m+1} + (1+\alpha)(-1)^{n+1}]
\} - \\
- \frac{4a\alpha\sinh a}{(m^2\pi^2+a^2)^2}
\{[\sinh a + a(-1)^{n+1}][(-1)^{m+n}-1]\} + \\
\left.+\tfrac{1}{2}\delta_{mn}+\alpha X_{mn}
-\tfrac{1}{2}(n^2\pi^2+a^2)\frac{\delta_{mn}}{a^2T}\right| = 0.
\end{multline}

A first approximation to the solution of equation~\eqref{eq:7-213} is obtained by
setting the $(1,1)$-element of the matrix equal to zero.  We find
\begin{equation}
\label{eq:7-214}
\tfrac{1}{2}(\pi^2+a^2)\frac{1}{Ta^2} = \tfrac{1}{2}a+\tfrac{1}{2} -
\frac{2a\pi^2(2+\alpha)}{(\pi^2+a^2)^2}
[(\sinh a\cosh a - a) + (\sinh a - a\cosh a)].
\end{equation}
On further simplification, this gives
\begin{equation}
\label{eq:7-215}
T = \frac{2}{2+\alpha}\,\frac{(\pi^2+a^2)^3}{a^2[1 -
16a^2\pi^2\cosh^2\!\frac{1}{2}a/(\pi^2+a^2)^2(\sinh a + a)]}.
\end{equation}
We observe that, apart from the factor $2/(2+\alpha)$, this expression for $T$
is identical with what was found in Chapter~II (\S~17, equation~(311)) for
the Rayleigh number for the simple B\'enard problem in the first
approximation for the case of two rigid boundaries.

Consequently, in this approximation (cf.\ equation~\eqref{eq:7-199}),
\begin{equation}
\label{eq:7-216}
T_c = \frac{2}{2+\alpha}\times 1715 = \frac{3430}{1+\mu}
\quad \text{and} \quad
a_{\min} = 3{\cdot}12.
\end{equation}

We shall see below (\S\,(b)) that for $0 < \mu < 1$ equation~\eqref{eq:7-216} gives
values for the critical Taylor number which do not differ from those
obtained in the higher approximations by more than one per cent.  The
reason for this relatively high accuracy of the solution in the first
approximation will be made apparent in \S\,(d).

\subsection*{(b) \emph{Numerical results}}
A method of solving the infinite order characteristic equation which~\eqref{eq:7-213} provides for $T$ would be to set the determinant formed by the first
$n$ rows and columns of the secular matrix equal to zero and let $n$ take
increasingly larger values.  In practice, the usefulness of this method will
depend largely on how rapidly the lowest positive root of the resulting
equation of order $n$ tends to its limit as $n \to \infty$.  It appears that for the
problem on hand, the process converges quite rapidly.

In Table~XXXII the values of $T$ obtained with the aid of equation~\eqref{eq:7-213} in the different approximations are listed for those values of $a$
(for the assigned $\mu$'s) at which it was found (by trial and error) that $T$
attained its minimum value.  From an examination of this table, it
would appear that for $\mu \geqslant -1{\cdot}0$, the third approximation provides $T$
to well within one per cent of the true value.  For $-3{\cdot}0 \leqslant \mu \leqslant -1{\cdot}0$,
the calculations were carried out to as high approximations as seemed

\begin{table}[htbp]
\centering
\caption{The critical Taylor numbers and the associated wave numbers for different
values of $\mu$}
\label{tab:XXXII}
\begin{tabular}{ccccccccc}
\hline
& & First & Second & Third & Fourth & Fifth & Sixth \\
$\mu$ & $a$ & approxi- & approxi- & approxi- & approxi- & approxi- & approxi- \\
& & mation & mation & mation & mation & mation & mation \\
\hline
$1{\cdot}00$ & $3{\cdot}12$ & $1715{\cdot}1$ & $1725{\cdot}2$ & $1708$ & & & \\
$0{\cdot}50$ & $3{\cdot}12$ & $1286{\cdot}7$ & & & & & \\
$0{\cdot}25$ & $3{\cdot}12$ & $2744{\cdot}1$ & $2736{\cdot}4$ & & & & \\
$-0{\cdot}25$ & $3{\cdot}12$ & $3430{\cdot}9$ & $3430{\cdot}6$ & $3399{\cdot}3$ & & & \\
$-0{\cdot}50$ & $3{\cdot}25$ & $4573{\cdot}4$ & $4472{\cdot}5$ & & & & \\
$-0{\cdot}50$ & $3{\cdot}25$ & $6800{\cdot}2$ & $6430{\cdot}6$ & $6417{\cdot}6$ & & $6148$ \\
$-0{\cdot}75$ & $3{\cdot}34$ & & $9422{\cdot}7$ & $9434{\cdot}6$ & & & \\
$-1{\cdot}00$ & $3{\cdot}50$ & $11514$ & $11840$ & $11500$ & & & \\
$-1{\cdot}25$ & $3{\cdot}75$ & & & $16847$ & & & \\
$-1{\cdot}50$ & $3{\cdot}75$ & & $18015$ & $18735$ & & $18077$ & \\
$-1{\cdot}75$ & $4{\cdot}00$ & & & $30497$ & $30458$ & & \\
$-2{\cdot}00$ & $4{\cdot}00$ & & & $45000$ & $40992$ & & \\
$-2{\cdot}50$ & $4{\cdot}00$ & & & $67189$ & & & \\
$-3{\cdot}00$ & & & & & $195810$ & $177107$ & $177108$ \\
$-3{\cdot}00$ & & & & & $330748$ & $301116$ & $302496$ \\
$-3{\cdot}09$ & & & & & & & $304452$ \\
\hline
\end{tabular}
\end{table}

necessary.  It would appear that for $\mu < -3{\cdot}0$, one should go to approximations higher than the sixth to achieve comparable accuracy.  Fortunately, the considerations to be set out in \S\,(c) below make it unnecessary for us to obtain the solutions for $\mu < -3{\cdot}0$ by the present method.

\figurePlaceholder{69}{The critical Taylor number $T_c$ for the onset of instability as a function
of $\mu = \Omega_2/\Omega_1$.}

\figurePlaceholder{70}{The wavelength of the disturbance (in units of the gap width, $d$)
manifested at onset of instability as a function of $\mu = \Omega_2/\Omega_1$.}

In Table~XXXIII the coefficients $\mathscr{C}_m$ ($= C_m/(m^2\pi^2+a^2)^{1/2}$) for the marginal state in the expansion~\eqref{eq:7-205} for $u$ are given.  This table also includes
the critical Taylor numbers and the associated wave numbers.  The
$(T_c, \mu)$ and the $[a(T_c), \mu]$-relationships are further illustrated in Figs.~69
and 70.

\begin{table}[htbp]
\centering
\caption{The critical Taylor numbers and related constants for various values of $\mu$}
\label{tab:XXXIII}
\small
\begin{tabular}{cccccccc}
\hline
$\mu$ & $a$ & $T_c$ & $\mathscr{C}_1/\mathscr{C}_1$ & $\mathscr{C}_2/\mathscr{C}_1$ & $\mathscr{C}_3/\mathscr{C}_1$ & $\mathscr{C}_4/\mathscr{C}_1$ & $\mathscr{C}_5/\mathscr{C}_1$ \\
\hline
$1{\cdot}00$ & $3{\cdot}12$ & $1{\cdot}208\times 10^3$ & $1$ & $0{\cdot}001324$ & & $0{\cdot}001147$ & \\
$0{\cdot}50$ & $3{\cdot}12$ & $2{\cdot}273\times 10^3$ & & $0{\cdot}001146$ & & & \\
$0{\cdot}25$ & $3{\cdot}12$ & $2{\cdot}900\times 10^3$ & & $0{\cdot}001304$ & & & \\
$0$ & $3{\cdot}14$ & $3{\cdot}4 \times 10^3$ & & $0{\cdot}000537$ & & $0{\cdot}000144$ & \\
$-0{\cdot}25$ & $3{\cdot}14$ & $4{\cdot}617\times 10^3$ & & $0{\cdot}001156$ & & $0{\cdot}000177$ & \\
$-0{\cdot}50$ & $3{\cdot}20$ & $6{\cdot}417\times 10^3$ & & $0{\cdot}001856$ & & & \\
$-0{\cdot}60$ & $3{\cdot}30$ & $7{\cdot}688\times 10^3$ & & $0{\cdot}003159$ & & $0{\cdot}001194$ & \\
$-0{\cdot}75$ & $3{\cdot}30$ & $1{\cdot}082\times 10^4$ & & $0{\cdot}003195$ & & & \\
$-0{\cdot}90$ & $3{\cdot}40$ & $1{\cdot}494\times 10^4$ & & $0{\cdot}04742$ & & & \\
$-1{\cdot}00$ & $3{\cdot}50$ & $1{\cdot}196\times 10^4$ & & $0{\cdot}00463$ & & & \\
$-1{\cdot}00$ & $4{\cdot}00$ & $1{\cdot}868\times 10^4$ & & $0{\cdot}02139$ & & $0{\cdot}000929$ & $-0{\cdot}000439$ \\
$-1{\cdot}25$ & $4{\cdot}00$ & & & $0{\cdot}04916$ & & $+0{\cdot}000421$ & \\
$-1{\cdot}25$ & $4{\cdot}06$ & $4{\cdot}019\times 10^4$ & & $+0{\cdot}02381$ & & $+0{\cdot}000116$ & $-0{\cdot}00012$ \\
$-1{\cdot}50$ & $4{\cdot}06$ & & & & & & \\
$-1{\cdot}75$ & $4{\cdot}06$ & $6{\cdot}528\times 10^4$ & & $+0{\cdot}04009$ & & & \\
$-2{\cdot}00$ & $4{\cdot}50$ & $9{\cdot}558\times 10^4$ & & $+0{\cdot}04396$ & & $-0{\cdot}001808$ & \\
$-2{\cdot}50$ & $5{\cdot}00$ & $1{\cdot}721\times 10^5$ & & $+0{\cdot}05804$ & & $+0{\cdot}001777$ & $-0{\cdot}000927$ \\
$-3{\cdot}00$ & & $3{\cdot}015\times 10^5$ & & $+0{\cdot}07499$ & & $+0{\cdot}001045$ & \\
$-3{\cdot}09$ & $5{\cdot}14$ & & & & & & $+0{\cdot}001045$ \\
\hline
\end{tabular}
\end{table}

Fig.~71\,$a$ shows the profile of the velocity $u$ for $\mu = -3{\cdot}0$; and in Fig.~71\,$b$ the corresponding cell pattern is exhibited.  For $\mu = -3{\cdot}0$, the

\figurePlaceholder{71}{The velocity profile (a) and the cell pattern (b) for $\mu=-3$ and
$a=8{\cdot}14$.  In (a) the velocity is normalized to unit maximum amplitude; in
(b) the line separating the positive and the negative values of the stream
function $\psi$ is indicated by the dashed line (note that this does not occur at
the nodal surface which is at $\zeta=0{\cdot}25$); (c) is the corresponding cell pattern
for the lowest inviscid mode for the same value of $a$ and $\mu$.}

nodal surface occurs at $\zeta = 0{\cdot}25$; and it is of interest to observe how
rapidly the amplitude of $u$ decreases as we go beyond this point.  The
velocity profile illustrated in Fig.~71 should be contrasted with that
in Fig.~67 which corresponds to the lowest inviscid mode for the same
value of $a$ and $\mu$.

\subsection*{(c) \emph{An alternative method of solution}}
Equations~\eqref{eq:7-201}--\eqref{eq:7-202} present the first example of a characteristic
value problem that we have encountered which is not self-adjoint.  For
this reason, it may be of interest to describe an alternative method of
solution to determine the principles which might guide one in devising
methods of solution which may be as equally convergent as the one we
have followed in \S\,(b).

Eliminating $u$ between equations~\eqref{eq:7-201} and~\eqref{eq:7-202}, we obtain
together with the boundary conditions
\begin{equation}
\label{eq:7-217}
(D^2-a^2)^3 v = -Ta^2(1+\alpha\zeta)v
\end{equation}
\begin{equation}
v = (D^2-a^2)v = D(D^2-a^2)v = 0
\quad \text{for } \zeta = 0 \text{ and } 1.
\label{eq:7-218}
\end{equation}
Rewriting equation~\eqref{eq:7-217} in the manner
\begin{equation}
\label{eq:7-219}
(D^2-a^2)^3 v = (1+\alpha\zeta)\psi
\end{equation}
and
\begin{equation}
\label{eq:7-220}
\psi = -Ta^2 v,
\end{equation}
we expand $\psi$ and $v$ in the forms
\begin{equation}
\label{eq:7-221}
\psi = \sum_{m=1}^{\infty} C_m \sin m\pi\zeta
\quad \text{and} \quad
v = \sum_{m=1}^{\infty} C_m\, v_m,
\end{equation}
where $v_m$ is the solution of the equation
\begin{equation}
\label{eq:7-222}
(D^2-a^2)^3 v_m = (1+\alpha\zeta)\sin m\pi\zeta
\end{equation}
which satisfies the boundary conditions~\eqref{eq:7-218}.  We then insert the
expansions~\eqref{eq:7-221} in~\eqref{eq:7-220} to obtain the secular equation.  The general solution of equation~\eqref{eq:7-222} is readily seen to be
\begin{multline}
\label{eq:7-223}
v_m = -\frac{1}{(m^2\pi^2+a^2)^2}
\biggl\{
A_1^{(m)}\cosh a\zeta + A_2^{(m)}\,\zeta\cosh a\zeta
+ A_1^{(m)}\,\zeta^2\cosh a\zeta + \\
+ B_1^{(m)}\sinh a\zeta
+ B_2^{(m)}\,\zeta\sinh a\zeta
+ B_2^{(m)}\,\zeta^2\sinh a\zeta + \\
+(1+\alpha\zeta)\sin m\pi\zeta
+ \frac{6\alpha m\pi}{m^2\pi^2+a^2}\cos m\pi\zeta
\biggr\},
\end{multline}
where $A_1^{(m)},\ldots, B_2^{(m)}$ are constants of integration to be determined by the
boundary conditions~\eqref{eq:7-218}.  These latter conditions lead to the equations:
\begin{align}
\label{eq:7-224}
A_1^{(m)} &= -\frac{6m\pi\alpha}{m^2\pi^2+a^2}, &
aA_1^{(m)}+3\,B_2^{(m)} &= \frac{m\pi}{2a}(m^2\pi^2+a^2),\nonumber\\
A_2^{(m)} + a\,B_2^{(m)} &= 2m\pi\alpha, & & \nonumber\\[3pt]
A_1^{(m)}\sinh a + A_2^{(m)}(&2a\sinh a + \cosh a) + \nonumber\\
+ B_1^{(m)}&a\cosh a + B_2^{(m)}(2a\cosh a + 3\sinh a)\,a\nonumber\\
&= (-1)^{m+1}\frac{6m\pi\alpha}{m^2\pi^2+a^2},\\[3pt]
A_1^{(m)}a\cosh a + A_2^{(m)}&(2a\cosh a + 3\sinh a)\,a + \nonumber\\
+ B_1^{(m)}&a\sinh a + B_2^{(m)}(2a\sinh a + 3\cosh a)\,a \nonumber\\
&= \frac{m\pi}{2a}(m^2\pi^2+a^2)(1+\alpha)(-1)^m.
\end{align}
On solving these equations, we find that
\begin{align}
\label{eq:7-207a}
A_1^{(m)} &= -\frac{4m\pi\alpha}{m^2\pi^2+a^2}, \nonumber\\
B_1^{(m)} &= \frac{m\pi}{m^2\pi^2+a^2}\,
\{\alpha+\beta_m(\sinh a + a\cosh a)-\gamma_m\sinh a\}, \nonumber\\
A_2^{(m)} &= -\frac{m\pi}{m^2\pi^2+a^2}
\{(\sinh^2 a + a)(\sinh a + a\cosh a)-\gamma_m \sinh a\}, \nonumber\\
B_2^{(m)} &= \frac{m\pi}{m^2\pi^2+a^2}
\{(\sinh a\cosh a - a)(1+\alpha)(-1)^{m+1} - \gamma_m(\alpha\cosh a - \alpha\cosh a)\},
\end{align}
where
$\Delta = \sinh^2 a - a^2$,
$\beta_m = \frac{4\alpha}{m^2\pi^2+a^2}\{(-1)^{m+1}\cosh a\}$,
and
$\gamma_m = (-1)^{m+1}(1+\alpha) + \frac{4\alpha}{m^2\pi^2+a^2}\, a\sinh a$.

Now substituting for $\psi$ and $v$ their expansions in equation~\eqref{eq:7-220}, we
obtain
\begin{equation}
\label{eq:7-225}
\frac{1}{Ta^2}\sum_{m=1}^{\infty} C_m \sin m\pi\zeta
= -\sum_{m=1}^{\infty} C_m\, v_m;
\end{equation}
and this clearly leads to the secular equation
\begin{equation}
\label{eq:7-226}
\left|\frac{1}{Ta^2}\,\delta_{mn}
+2\langle n|m\rangle\right| = 0,
\end{equation}
where $\langle n|m\rangle$ stands for the matrix element
\begin{equation}
\label{eq:7-227}
\langle n|m\rangle = \int_0^1 v_m \sin n\pi\zeta\,d\zeta.
\end{equation}
(Observe that the matrix $\langle k|x|j\rangle$ is not Hermitian.)

On evaluating the matrix element $\langle n|m\rangle$, we find after some lengthy but
elementary reductions that the secular equation takes the explicit form
\begin{multline}
\label{eq:7-228}
\left|
\frac{nmm'\pi^2}{n^2\pi^2+a^2}\left(\frac{6\alpha}{m^2\pi^2+a^2}
+ \frac{4\alpha}{n^2\pi^2+a^2}\right)
\{(-1)^{m+n}-1\} + \right. \\
+ \frac{8a^2}{(m^2\pi^2+a^2)^2}\left\{
\frac{4\alpha^2}{4a^2}\right\}
\{1+(-1)^{m+n}+(-1)^m\}
\cosh a\{(-1)^{m+n}+(-1)^n\} + \\
+\alpha(-1)^m\,
\frac{1+(-1)^{m+n}\cosh a}{\sinh a + (-1)^{n+1}a}
+\alpha\left\{(-1)^{m+n+1}\frac{\sinh a}{\sinh a + (-1)^{n+1}a}\right\} + \\
\left.+\tfrac{1}{2}\delta_{mn}+\alpha X_{mn}
-\tfrac{1}{2}(n^2\pi^2+a^2)^2\frac{\delta_{mn}}{a^2 T}\right| = 0,
\end{multline}
where
\begin{equation}
\label{eq:7-229}
X_{mn} = \begin{cases}
0 & \text{if } m+n \text{ is even and } m \neq n,\\
\frac{1}{2} & \text{if } m = n,\\
\dfrac{4mn}{n^2-m^2}\left(\dfrac{3}{m^2\pi^2+a^2}
- \dfrac{1}{\pi(n^2-m^2)}\right) & \text{if } m+n \text{ is odd.}
\end{cases}
\end{equation}

One might have thought that, since equation~\eqref{eq:7-228} is based on the
solution of an equation of order six, the characteristic values derived
from the present equation in the different approximations would show
an even more rapid convergence than has been found by the method of
\S\,(a).  However, this is not the case: the results given by equations~\eqref{eq:7-213}
and~\eqref{eq:7-228} are practically indistinguishable.  Thus:
\begin{equation}
\label{eq:7-230}
\left.
\begin{aligned}
T &= 6417{\cdot}8 &&\text{(from equation~\eqref{eq:7-228} in 3rd approximation)}\\
  &= 6417{\cdot}1 &&\text{(from equation~\eqref{eq:7-213} in 3rd approximation)}
\end{aligned}
\right\}\\
\text{for } \mu = -0{\cdot}5 \text{ and } a = 3{\cdot}20
\end{equation}
and
\begin{equation*}
T\left\{
\begin{aligned}
&= 95624 &&\text{(from equation~\eqref{eq:7-228} in 4th approximation)}\\
&= 95625 &&\text{(from equation~\eqref{eq:7-213} in 4th approximation)}
\end{aligned}
\right\}\\
\end{equation*}
for $\mu = -2{\cdot}0$ and $a = 6{\cdot}05$.

The explanation for this near efficacy of the two methods probably
derives from the fact that the advantage gained in the present method
by solving an equation of order six has been lost in the artificial treatment
of the equation and the boundary conditions.  In the method described
in \S\,(a), the variables used had direct physical meanings, and the
boundary conditions relevant to each and the relations between them
were satisfied identically; and this is perhaps the real key to the problem.

\subsection*{(d) \emph{An approximate solution for $\mu \to 1$}}
From the results given in Table~XXXIII, it appears that the formula
\begin{equation}
\label{eq:7-231}
T_c = \frac{3416}{1+\mu}
\end{equation}
gives a very good approximation to the true values for $0 \leqslant \mu \leqslant 1$;
and, moreover, the wave number ($\sim 3{\cdot}12$) at which instability sets in,
hardly seems to depend on $\mu$ in the same range.  Formula~\eqref{eq:7-231} was
derived by Taylor by a method of solution very different from the ones
we have described; and we have also seen how effectively the same
formula is given by the \emph{first} approximation to the solution of the secular
equation~\eqref{eq:7-213}.

We shall now show how one can derive~\eqref{eq:7-231} by a simple perturbation
procedure applied to the solution~\eqref{eq:7-217}; at the same time, we
shall obtain a correction term to~\eqref{eq:7-231} which extends its range of validity.

It is convenient for our present purposes to translate the origin of the
coordinate system to be midway between the two cylinders.  Thus, let
\begin{equation}
\label{eq:7-232}
\zeta' = x + \tfrac{1}{2}.
\end{equation}
so that the limits of $x$ are $\pm\frac{1}{2}$.  In terms of $x$, equation~\eqref{eq:7-217} becomes
\begin{equation}
\label{eq:7-233}
(D^2-a^2)^3 v = -\lambda(1+\epsilon x)v,
\end{equation}
where
\begin{equation}
\label{eq:7-234}
\lambda = \tfrac{1}{2}Ta^2(1+\mu) \quad \text{and} \quad
\epsilon = -2\,\frac{1-\mu}{1+\mu},
\end{equation}
and the boundary conditions are
\begin{equation}
\label{eq:7-235}
u = (D^2-a^2)v = D(D^2-a^2)v = 0
\quad \text{for } x = \pm\tfrac{1}{2}.
\end{equation}
When $\epsilon = 0$, the characteristic value problem presented by equations~\eqref{eq:7-233} and~\eqref{eq:7-235} reduces to the one for the simple B\'enard problem (for
the case of two rigid boundaries) when formulated in terms of the
amplitude $\Theta$ of the fluctuations in the temperature (see Chapter~II,
\S\S~12 and 13\,(b)).  Let $\lambda_j = R_j a^2$ in the notation of Chapter~II for
$j = 0, 1, \ldots$, denote the different characteristic values of the equation
\begin{equation}
\label{eq:7-236}
(D^2-a^2)^3\Theta = -\lambda\Theta,
\end{equation}
for the boundary conditions
\begin{equation}
\label{eq:7-237}
\Theta = (D^2-a^2)\Theta = D(D^2-a^2)\Theta = 0
\quad \text{for } x = \pm\tfrac{1}{2}.
\end{equation}
and let the functions belonging to $\lambda_j$ be distinguished by a subscript $j$.
We know from Chapter~II, \S~13\,(b) that the functions\footnote{As defined here $\Theta$ and $W$ are the same functions which describe the perturbations
in the temperature and in the normal component of the velocity for the B\'enard problem
(see equations II~(90)).}
\begin{equation}
\label{eq:7-238}
\Theta_j \quad \text{and} \quad W_j = (D^2-a^2)\Theta_j
\end{equation}
are mutually orthogonal\footnote{Equation~II~(161) directly implies that
\[
\int \{(D\Theta_j)(D\Theta_k)+a^2\Theta_j\Theta_k\}\,dx = 0 \quad \text{if } j \neq k.
\]
Integrating by parts the first term in the integrand, we obtain
\[
\int \Theta_j(D^2-a^2)\Theta_k\,dx = \int \Theta_j\,W_k\,dx = 0 \quad (j \neq k).
\]}
when $j \neq k$; without loss of generality, we
may suppose that they are so normalized that
\begin{equation}
\label{eq:7-239}
\int_{-1/2}^{+1/2}\Theta_j\, W_k\, dx = \delta_{jk}.
\end{equation}

Turning to the solution of~\eqref{eq:7-233}, we first observe that according to
the results given in Table~III (Chapter~II) the lowest value of $\lambda/a^2$,
when $\epsilon = 0$, is 1708 and occurs for $a = 3{\cdot}117$.  By~\eqref{eq:7-234} the corresponding value of $T$ is $3416/(1+\mu)$; and this is exactly the same as~\eqref{eq:7-231}.
Thus, formula~\eqref{eq:7-231} appears as the zero-order term in a perturbation
series.  To obtain the higher-order terms, we expand the various quantities in powers of $\epsilon$.  Thus, we write
\begin{equation}
\label{eq:7-240}
\left.\begin{aligned}
v &= \Theta_0+\epsilon v^{(1)}+\epsilon^2 v^{(2)}+\cdots\\
\lambda &= \lambda_0+\epsilon\lambda^{(1)}+\epsilon^2\lambda^{(2)}+\cdots
\end{aligned}\right\}.
\end{equation}
Substituting these expansions in equation~\eqref{eq:7-233} and equating the
different powers of $\epsilon$, we obtain:
\begin{align}
\label{eq:7-241}
(D^2-a^2)^3 \Theta_0 &= -\lambda_0\Theta_0, \\
\label{eq:7-242}
(D^2-a^2)^3 v^{(1)} &= -(\lambda_0 v^{(1)}+\lambda_0 x\Theta_0+\lambda^{(1)}\Theta_0), \\
\label{eq:7-243}
(D^2-a^2)^3 v^{(2)} &= -(\lambda_0 v^{(2)}+\lambda_0 xv^{(1)}+\lambda^{(1)} v^{(1)}
+\lambda^{(2)}\Theta_0+\lambda^{(1)}x\Theta_0),
\end{align}
etc.

Equation~\eqref{eq:7-241} is clearly satisfied.  To solve equation~\eqref{eq:7-242}, we assume
for $v^{(1)}$ a solution of the form
\begin{equation}
\label{eq:7-244}
v^{(1)} = \sum_{j=0}^{\infty} A_j^{(1)}\Theta_j.
\end{equation}
Inserting this solution in~\eqref{eq:7-242}, we obtain
\begin{equation}
\label{eq:7-245}
\sum_{j=0}^{\infty} A_j^{(1)}\lambda_j\,\Theta_j
= \lambda_0 \sum_{j=0}^{\infty} A_j^{(1)}\Theta_j
+ \lambda_0 x\Theta_0 + \lambda^{(1)}\Theta_0.
\end{equation}
Multiply this equation by $W_k$ and integrate over the range of $x$.  Making
use of the orthogonality relation~\eqref{eq:7-239}, we obtain
\begin{equation}
\label{eq:7-246}
A_k^{(1)}\lambda_k = \lambda_0\, A_k^{(1)}
+ \lambda_0\langle k|x|0\rangle + \lambda^{(1)}\delta_{k0},
\end{equation}
where the notation
\begin{equation}
\label{eq:7-247}
\langle k|x|j\rangle = \int_{-1/2}^{+1/2} W_k\, x\,\Theta_j\, dx
\end{equation}
has been adopted.  (Observe that the matrix $\langle k|x|j\rangle$ is not Hermitian.)
For $k = 0$, equation~\eqref{eq:7-246} gives
\begin{equation}
\label{eq:7-248}
\lambda^{(1)} = -\lambda_0\langle 0|x|0\rangle = -\lambda_0
\int_{-1/2}^{+1/2} W_0\, x\, \Theta_0\, dx.
\end{equation}
Since the proper functions $W_0$ and $\Theta_0$ belonging to $\lambda_0$ are even, the matrix
element $\langle 0|x|0\rangle$ is zero; and we conclude
\begin{equation}
\label{eq:7-249}
\lambda^{(1)} = 0.
\end{equation}
Therefore, formula~\eqref{eq:7-231} is in error by a term which is of the \emph{second
order} in $\epsilon = -2(1-\mu)/(1+\mu)$; and this explains its success in representing the results of more exact calculations as well as it does.

For $k \neq 0$, equation~\eqref{eq:7-246} gives
\begin{equation}
\label{eq:7-250}
A_k^{(1)} = \frac{\lambda_0}{\lambda_k-\lambda_0}\,\langle k|x|0\rangle.
\end{equation}
The coefficient $A_0^{(1)}$ is left unspecified; but as we shall presently see, this
does not lead to any ambiguities in the solution for the physical problem.

Consider next equation~\eqref{eq:7-243}.  Since $\lambda^{(1)} = 0$, this equation reduces to
\begin{equation}
\label{eq:7-251}
(D^2-a^2)^3 v^{(2)} = -(\lambda_0 v^{(2)}
+ \lambda_0 xv^{(1)} + \lambda^{(1)}xv^{(1)} + \lambda^{(2)}\Theta_0).
\end{equation}
In solving this equation, we again expand $v^{(2)}$ in terms of $\Theta_j$.  Thus, with
the assumption
\begin{equation}
\label{eq:7-252}
v^{(2)} = \sum_{j=0}^{\infty} A_j^{(2)}\Theta_j,
\end{equation}
equation~\eqref{eq:7-251} gives
\begin{equation}
\label{eq:7-253}
\sum_{j=0}^{\infty} A_j^{(2)}\lambda_j\,\Theta_j
= \lambda_0\sum_{j=0}^{\infty} A_j^{(2)}\Theta_j
+ \lambda_0 xv^{(1)} + \lambda^{(2)}\Theta_0.
\end{equation}
Multiplying this equation by $W_k$ and integrating over the range of $x$, we
obtain
\begin{equation}
\label{eq:7-254}
A_k^{(2)}\lambda_k = \lambda_0\, A_k^{(2)}
+ \lambda_0 \int W_k\, xv^{(1)}\, dx + \lambda^{(2)}\delta_{k0}.
\end{equation}
Hence
\begin{equation}
\label{eq:7-255}
\lambda^{(2)} = -\lambda_0 \int_{-1/2}^{+1/2} W_0\, xv^{(1)}\, dx.
\end{equation}
Substituting for $v^{(1)}$ from equation~\eqref{eq:7-244}, we obtain
\begin{equation}
\label{eq:7-256}
\lambda^{(2)} = -\lambda_0 \sum_{j=0}^{\infty} A_j^{(1)}
\int_{-1/2}^{+1/2} W_0\, x\,\Theta_j\, dx.
\end{equation}
With the notation~\eqref{eq:7-247} for the matrix elements of $x$, we can write
\begin{equation}
\label{eq:7-257}
\lambda^{(2)} = -\lambda_0 \sum_{j=0}^{\infty} A_j^{(1)}\langle 0|x|j\rangle.
\end{equation}
Since $\langle 0|x|0\rangle = 0$, the term $j = 0$ in~\eqref{eq:7-257} makes no contribution; and
substituting for $A_j^{(1)}$ ($j \neq 0$) from equation~\eqref{eq:7-250}, we have
\begin{equation}
\label{eq:7-258}
\lambda^{(2)} = -\lambda_0^2\sum_{j \neq 0}
\frac{\langle 0|x|j\rangle\langle j|x|0\rangle}{\lambda_j-\lambda_0}.
\end{equation}
This expression for the second-order change in the characteristic
value is very similar to the expression one has in the quantum theory for the
second-order change in the energy levels of an atomic system caused by
a perturbation.  An important difference, however, is the non-Hermitian
character of the matrix representing the perturbation.

In the particular problem to which equation~\eqref{eq:7-258} refers, the `energy
levels' are so widely spaced (e.g.\ $\lambda_1 \sim 15\lambda_0$) that it will be sufficient to
retain only the first term in the infinite sum representing $\lambda^{(2)}$.  Thus, we
may write
\begin{equation}
\label{eq:7-259}
\lambda^{(2)} \simeq -\frac{\lambda_0^2}{\lambda_1-\lambda_0}
\langle 0|x|1\rangle\langle 1|x|0\rangle,
\end{equation}
To the second order in $\epsilon$, the desired expression for the characteristic
value is
\begin{equation}
\label{eq:7-260}
\lambda = \lambda_0\!\left\{1
- \epsilon^2\frac{\lambda_0}{\lambda_1-\lambda_0}
\langle 0|x|1\rangle\langle 1|x|0\rangle\right\}.
\end{equation}
Substituting for $\lambda$ from equation~\eqref{eq:7-234} and remembering that $\lambda_j = R_j\, a^2$,
we can rewrite equation~\eqref{eq:7-260} in the form
\begin{equation}
\label{eq:7-261}
T = \frac{2R_0}{1+\mu}\!\left\{1
- \epsilon^2\frac{R_0}{R_1-R_0}\langle 0|x|1\rangle\langle 1|x|0\rangle\right\},
\end{equation}
where $R_0$ and $R_1$ are the Rayleigh numbers for the lowest even and odd
modes for the assigned $a^2$.

To obtain the critical Taylor number for the onset of instability, we
must minimize the quantity on the right-hand side of equation~\eqref{eq:7-261}
as a function of $a$.  The value of $a$ at which $R_0$ ($= 1708$, by $a_{\min}^{(0)} = 3{\cdot}117$) at which $R_0$ attains its
minimum value, $R_c$ ($= 1708$), by an amount of order $\epsilon^2$.  At the displaced position of the new minimum, $R_0$ will differ from $R_c$ only by
an amount of order $\epsilon^4$.  Therefore, to order $\epsilon^2$, we may write
\begin{equation}
\label{eq:7-262}
T_c = \frac{2R_c}{1+\mu}\!\left\{1
- 4\!\left(\frac{1-\mu}{1+\mu}\right)^{\!2}
\frac{R_0}{R_1^*-R_c}\,
\langle 0|x|1\rangle\langle 1|x|0\rangle\right\},
\end{equation}
where $R_1^*$ ($= 2{\cdot}4982\times 10^4$) is the Rayleigh number for the first odd
mode for $a_{\min}^{(0)}$ (see Table~I, p.~38).  [In~\eqref{eq:7-262} we have substituted for $\epsilon$ its value
$-2(1-\mu)/(1+\mu)$.]

The solution for the first odd mode which is needed for the evaluation
of the matrix elements in~\eqref{eq:7-262} can be found by the method described in
Chapter~II, \S~15\,(b).  And it is found that the numerical form of equation~\eqref{eq:7-262} is
\begin{equation}
\label{eq:7-263}
T_c = \frac{3416}{1+\mu}
\!\left\{1-7{\cdot}61\times 10^{-2}
\!\left(\frac{1-\mu}{1+\mu}\right)^{\!2}\right\}.
\end{equation}
This formula gives $T_c = 3{\cdot}390\times 10^3$ and $6{\cdot}36\times 10^3$ for $\mu = 0$ and $-0{\cdot}5$,
respectively; these values should be compared with $3{\cdot}390\times 10^3$ and
$6{\cdot}42\times 10^3$ given by the `exact' calculations.

\subsection*{(e) \emph{The asymptotic behaviour for $(1-\mu) \to \infty$}}
From an examination of the results given in Table~XXXIII, it appears
that for $(1-\mu) \to \infty$ the following asymptotic relations obtain:
\begin{equation}
\label{eq:7-264}
T_c \to \tau(1-\mu)^4 \quad \text{and} \quad
a(T_c) \to q(1-\mu) \quad \text{as } (1-\mu) \to \infty,
\end{equation}
where $\tau$ and $q$ are certain constants.  Thus, from the tabulated values
we find:
\begin{equation}
\label{eq:7-265}
\frac{T_c}{(1-\mu)^4}\left\{
\begin{aligned}
&= 1180{\cdot}5\\
&= 1180{\cdot}2\\
&= 1181{\cdot}6
\end{aligned}
\right.
\quad \text{and} \quad
\frac{a(T_c)}{1-\mu}\left\{
\begin{aligned}
&= 2{\cdot}03\\
&= 2{\cdot}03\\
&= 2{\cdot}035
\end{aligned}
\right.
\quad \text{for } 1-\mu\left\{
\begin{aligned}
&= 3{\cdot}0\\
&= 3{\cdot}5\\
&= 4{\cdot}0.
\end{aligned}
\right.
\end{equation}
The fact that asymptotic relations such as~\eqref{eq:7-264} should obtain can be
readily understood.  For, it will be recalled that the origin of the instability, in the case when the two cylinders are in counter rotation, lies
in the violation of the Rayleigh criterion in the layers adjoining the inner
cylinder and extending not much beyond the nodal surface.  On these
grounds, one may expect that as $(1-\mu) \to \infty$ the critical Taylor
number for the onset of instability must be comparable to the critical Taylor
number for the case $\mu = 0$, and for a gap width equal to the distance
$d_0$ of the nodal surface to the inner cylinder.  For the distribution of the
angular velocity given by~\eqref{eq:7-191},
\begin{equation}
\label{eq:7-266}
d_0 = d/(1-\mu).
\end{equation}
Since the gap width occurs with the fourth power in the definition of $T$,
the foregoing arguments will lead one to expect an asymptotic behaviour
of the type
\begin{equation}
\label{eq:7-267}
T_c \to \tau(1-\mu)^4,
\end{equation}
where $\tau$ is comparable to the Taylor number for $\mu = 0$.  For the same
reasons, one may also expect that the wave number of the associated
disturbance will show the behaviour
\begin{equation}
\label{eq:7-268}
a(T_c) \to q(1-\mu),
\end{equation}
where $q$ is comparable to the wave number appropriate for $\mu = 0$.

The fact that the constants $\tau$ ($\sim 1182$) and $q$ ($\sim 2{\cdot}035$) indicated by
the calculations are substantially different from the values 3416 and
$3{\cdot}12$ appropriate for $\mu = 0$ should not cause much surprise; for, at
$\zeta = d_0/d$, the conditions which prevail are far from those at $\zeta = 1$ for
$\mu = 0$; and it is also not true, even asymptotically, that the disturbance
does not extend beyond the nodal surface.

While the foregoing arguments give physical grounds for expecting
asymptotic relations of the forms~\eqref{eq:7-267} and~\eqref{eq:7-268}, we can equally trace
their origin to the analytical structure of the underlying characteristic
value problem.  Thus, with the substitution
\begin{equation}
\label{eq:7-269}
x = a\zeta,
\end{equation}
equation~\eqref{eq:7-217} becomes
\begin{equation}
\label{eq:7-270}
(D^2-1)^3 v = -\frac{T}{a^4}\!\left(1-\frac{1-\mu}{a}\,x\right)v;
\end{equation}
and the corresponding boundary conditions are
\begin{equation}
\label{eq:7-271}
v = (D^2-1)v = D(D^2-1)v = 0
\quad \text{for } x = 0 \text{ and } a.
\end{equation}
Suppose now that
\begin{equation}
\label{eq:7-272}
\frac{T}{a^4} \to \lambda \quad \text{and} \quad
\frac{1-\mu}{a} \to \gamma \quad
\text{as } (1-\mu) \to \infty \text{ and } a \to \infty.
\end{equation}
In this limit equations~\eqref{eq:7-270} and~\eqref{eq:7-271} become
\begin{equation}
\label{eq:7-273}
(D^2-1)^3 v = -\lambda(1-\gamma x)v
\end{equation}
and
\begin{equation}
\label{eq:7-274}
v = (D^2-1)v = D(D^2-1)v = 0
\quad \text{for } x = 0 \text{ and } \infty.
\end{equation}
For an assigned value of $\gamma$, equations~\eqref{eq:7-273} and~\eqref{eq:7-274} will determine
a sequence of possible values of $\lambda$; and to determine the constants $\tau$
and $q$ in the asymptotic relations~\eqref{eq:7-267} and~\eqref{eq:7-268}, we must find the
minimum of $\lambda/\gamma^4$ as a function of $\gamma$.  Thus,
\begin{equation}
\label{eq:7-275}
\tau = \min_\gamma\!\left\{\frac{\lambda}{\gamma^4}\right\},
\end{equation}
and $q$ is the reciprocal of the value of $\gamma$ at which $\lambda/\gamma^4$ attains its minimum
value.

The characteristic value problem presented by equations~\eqref{eq:7-273} and~\eqref{eq:7-274} has not been solved satisfactorily at the present time: for various
reasons, the methods which have proved successful in other cases seem
to fail when applied to the present problem.
